\documentclass[UTF8, 12pt, fontset=fandol, oneside]{ctexart}
\usepackage{researchnotes}

\title{第二章\quad 矩阵}
\author{}
\date{}

\begin{document}

\maketitle

\section*{§2.1\quad 矩阵的概念}

矩阵是什么？它有什么用处？

在上一章我们学习了 $n$ 阶行列式的概念. 一个 $n$ 阶行列式从形式上看无非是将 $n^2$ 个元素排成 $n$ 行与 $n$ 列：
\[
\begin{array}{cccc}
a_{11} & a_{12} & \cdots & a_{1n}\\
a_{21} & a_{22} & \cdots & a_{2n}\\
\vdots & \vdots & & \vdots\\
a_{n1} & a_{n2} & \cdots & a_{nn}
\end{array},
\]
在许多实际问题中，还会碰到由若干个数排成行与列的长方形数组. 在研究问题时常常需要把它作为一个整体来处理，这就需要我们引进矩阵的概念.

\noindent\textbf{定义 2.1.1}\quad 由 $mn$ 个数 $a_{ij}(i=1,2,\cdots,m;\ j=1,2,\cdots,n)$ 排成 $m$ 行、$n$ 列的矩形阵列：
\[
\begin{array}{cccc}
a_{11} & a_{12} & \cdots & a_{1n}\\
a_{21} & a_{22} & \cdots & a_{2n}\\
\vdots & \vdots & & \vdots\\
a_{m1} & a_{m2} & \cdots & a_{mn}
\end{array}
\]
称为 $m$ 行 $n$ 列矩阵，简称为 $m\times n$ 矩阵 (或 $m\times n$ 阵).

矩阵常用大写英文字母来表示，且为了写得紧凑便于分辨，往往用括弧将上述矩阵括起来. 比如上面定义中的矩阵可记为
\[
A=
\left(\begin{array}{cccc}
a_{11} & a_{12} & \cdots & a_{1n}\\
a_{21} & a_{22} & \cdots & a_{2n}\\
\vdots & \vdots & & \vdots\\
a_{m1} & a_{m2} & \cdots & a_{mn}
\end{array}\right).
\]
有时为了简单起见，就记为 $A=(a_{ij})_{m\times n}$. 这里 $m$ 写在前面表示共有 $m$ 行，$n$ 写在后面表示共有 $n$ 列.

$a_{ij}(i=1,2,\cdots,m;\ j=1,2,\cdots,n)$ 称为矩阵 $A$ 的元素，前一个足标 $i$ 表示这个元素在第 $i$ 行，后一个足标 $j$ 表示这个元素在第 $j$ 列. $a_{ij}$ 称为矩阵 $A$ 的第 $i$ 行第 $j$ 列元素，或简称为 $A$ 的第 $(i,j)$ 元素. 矩阵的元素通常用英文小写字母表示或直接用数字表示.

如果矩阵 $A$ 的元素全是实数，则称 $A$ 为实矩阵；如果 $A$ 的元素全为复数，则称之为复矩阵. 所有元素均为零的矩阵，叫零矩阵，记为 $O$. 但是必须注意这个 $O$ 不是一个数，而是一个矩阵. 有时为了强调这是一个 $m\times n$ 阵，可写为 $O_{m\times n}$.

若矩阵的行列数相等，则称之为方阵. 含有 $n$ 行及 $n$ 列的矩阵称为 $n$ 阶方阵 (亦称为 $n$ 阶矩阵)，行列数不相等的矩阵称为长方阵. 方阵在矩阵理论中占有特别重要的位置. 若 $A=(a_{ij})$ 是 $n$ 阶方阵，则元素 $a_{11},a_{22},\cdots,a_{nn}$ 称为 $A$ 的主对角线. 若一个方阵除了主对角线上的元素外其余元素都等于零，就称之为对角阵. 对角阵的形状为
\[
A=
\left(\begin{array}{cccc}
a_{11} & 0 & \cdots & 0\\
0 & a_{22} & \cdots & 0\\
\vdots & \vdots & & \vdots\\
0 & 0 & \cdots & a_{nn}
\end{array}\right).
\]
上述对角阵可简记为 $\operatorname{diag}\{a_{11},a_{22},\cdots,a_{nn}\}$. 若进一步有 $a_{11}=a_{22}=\cdots=a_{nn}=1$，则称这个矩阵为单位阵. $n$ 阶单位阵通常记为 $I_n$：
\[
I_n=
\left(\begin{array}{cccc}
1 & 0 & \cdots & 0\\
0 & 1 & \cdots & 0\\
\vdots & \vdots & & \vdots\\
0 & 0 & \cdots & 1
\end{array}\right).
\]

一个 $n$ 阶方阵，如果它的主对角线以下的元素都等于零，即它具有下列形状：
\[
A=
\left(\begin{array}{cccc}
a_{11} & a_{12} & \cdots & a_{1n}\\
0 & a_{22} & \cdots & a_{2n}\\
\vdots & \vdots & & \vdots\\
0 & 0 & \cdots & a_{nn}
\end{array}\right),
\]
则称 $A$ 为上三角阵. 同样地，若 $A$ 的主对角线上面的元素全为零，则称 $A$ 为下三角阵. 下三角阵的形状为
\[
A=
\left(\begin{array}{cccc}
a_{11} & 0 & \cdots & 0\\
a_{21} & a_{22} & \cdots & 0\\
\vdots & \vdots & & \vdots\\
a_{n1} & a_{n2} & \cdots & a_{nn}
\end{array}\right).
\]

我们定义了矩阵的概念. 从定义可以看出，矩阵可以是各种各样的. 为了判断其异同，我们必须说明什么叫两个矩阵相等. 简单地说，两个矩阵相等要求它们``完全一样''. 即若 $A=(a_{ij})_{m\times n}, B=(b_{ij})_{s\times t}$，则 $A=B$ 当且仅当 $m=s,n=t$，且 $a_{ij}=b_{ij}$ 对所有 $i,j$ 都成立. 换言之，两个矩阵相等要求它们含有相同数目的行及相同数目的列，且第 $(i,j)$ 元素都对应相等.

矩阵相等的概念并不复杂，但有时也会产生混淆. 比如下面两个矩阵：
\[
\left(\begin{array}{cc}
0 & 0\\
0 & 0
\end{array}\right),\quad
\left(\begin{array}{ccc}
0 & 0 & 0\\
0 & 0 & 0
\end{array}\right)
\]
都是零矩阵且常用同一个符号表示，但它们不相等. 因为一个是 $2\times 2$ 矩阵，另一个是 $2\times 3$ 矩阵. 下面两个矩阵也不相等，虽然它们都是二阶矩阵且都含有两个 $0$ 与两个 $1$：
\[
\left(\begin{array}{cc}
1 & 1\\
0 & 0
\end{array}\right),\quad
\left(\begin{array}{cc}
1 & 0\\
0 & 1
\end{array}\right).
\]

矩阵的每一行称为这个矩阵的一个行向量，同样它的每一列称为它的一个列向量. 一个 $1\times n$ 矩阵：
\[
(a_1,a_2,\cdots,a_n)
\]
称为 $n$ 维行向量；一个 $n\times 1$ 矩阵：
\[
\left(\begin{array}{c}
a_1\\
a_2\\
\vdots\\
a_n
\end{array}\right)
\]
称为 $n$ 维列向量.

读者可能会问，矩阵和行列式有什么不同？从上面的定义我们可以看出，矩阵就是矩形数组，而行列式 (注意仅对方阵而言) 是指这个方阵所对应的一个数值. 若 $A$ 是 $n$ 阶方阵，则我们用 $|A|$ 或 $\det A$ 表示矩阵 $A$ 的行列式. 注意对长方阵而言，谈论其行列式显然没有意义.

\section*{§2.2\quad 矩阵的运算}

初等代数研究的是数，但是数不仅仅是数字的集合，它们之间还可以进行运算. 因此，初等代数可以说是研究数及其运算的学科. 现在我们要研究矩阵，即数组. 数组之间也可以定义运算，研究矩阵及其运算规律是高等代数的基本任务. 读者要注意，矩阵的代数运算是应实际需要引进的而不是数学家的随意创造. 在这一节里我们将介绍矩阵的加减法、数乘、乘法 (包括乘方)、转置与共轭. 这些运算有些与通常数字的运算相似，有些则有很大的差别. 读者务须熟练掌握这些概念.

\subsection*{一、矩阵的加减法}

\noindent\textbf{定义 2.2.1}\quad 设有两个 $m\times n$ 矩阵 $A=(a_{ij}), B=(b_{ij})$，定义 $A+B$ 是一个 $m\times n$ 矩阵且 $A+B$ 的第 $(i,j)$ 元素等于 $a_{ij}+b_{ij}$，即
\[
A+B=(a_{ij}+b_{ij}).
\]
例如
\[
\left(\begin{array}{ccc}
1 & 2 & 0\\
-1 & 3 & 1
\end{array}\right)
+
\left(\begin{array}{ccc}
1 & -2 & 1\\
-1 & -3 & 0
\end{array}\right)
=
\left(\begin{array}{ccc}
2 & 0 & 1\\
-2 & 0 & 1
\end{array}\right).
\]
在进行矩阵加法时需特别注意，相加的两个矩阵的行数与列数必须分别对应相等. 凡不满足这个条件的矩阵是不可以相加的.

一个 $m\times n$ 矩阵 $A$ 与一个 $m\times n$ 零矩阵相加显然仍等于 $A$，即
\[
A+O=A.
\]
因此零矩阵在矩阵加法中的作用与数字 $0$ 在数的加法中的作用类似.

矩阵的减法可以看成是矩阵加法的逆运算，因此行数与列数分别相等的两个矩阵才可以相减. 若 $A=(a_{ij})_{m\times n}, B=(b_{ij})_{m\times n}$，则
\[
A-B=(a_{ij}-b_{ij}).
\]
例如
\[
\left(\begin{array}{cc}
1 & 3\\
5 & 2\\
-1 & 0
\end{array}\right)
-
\left(\begin{array}{cc}
1 & 1\\
3 & 0\\
0 & 1
\end{array}\right)
=
\left(\begin{array}{cc}
0 & 2\\
2 & 2\\
-1 & -1
\end{array}\right).
\]

两个相等的矩阵相减为零矩阵，即若 $A=B$，则
\[
A-B=O.
\]

我们还可以定义负矩阵，设 $A=(a_{ij})$，定义 $-A=(-a_{ij})$. 显然
\[
A+(-A)=O.
\]

\noindent\textbf{矩阵加减法运算规则}\quad 矩阵的加减法运算适合下列规则：

(1) 交换律：$A+B=B+A$;

(2) 结合律：$(A+B)+C=A+(B+C)$;

(3) $O+A=A+O=A$;

(4) $A+(-B)=A-B$.

这些规则的验证是十分容易的，请读者自己完成.

\subsection*{二、矩阵的数乘}

\noindent\textbf{定义 2.2.2}\quad 设 $A$ 是一个 $m\times n$ 矩阵，$A=(a_{ij})_{m\times n}$，$c$ 是一个数，定义 $cA=(ca_{ij})_{m\times n}$. $cA$ 称为数 $c$ 与矩阵 $A$ 的数乘.

上述定义告诉我们：一个数与一个矩阵的数乘等于将这个数乘以矩阵的每一项 (即每个元素) 所得的矩阵. 例如：
\[
3\cdot
\left(\begin{array}{cc}
1 & 1\\
3 & 0\\
0 & 1
\end{array}\right)
=
\left(\begin{array}{cc}
3 & 3\\
9 & 0\\
0 & 3
\end{array}\right).
\]
矩阵 $A$ 的负矩阵也可以看成是 $-1$ 与 $A$ 的数乘.

\noindent\textbf{矩阵数乘运算规则}\quad 矩阵的数乘运算适合下列规则：

(1) $c(A+B)=cA+cB$;

(2) $(c+d)A=cA+dA$;

(3) $(cd)A=c(dA)$;

(4) $1\cdot A=A$;

(5) $0\cdot A=O$.

\subsection*{三、矩阵的乘法}

下面要定义的矩阵乘法是矩阵运算中最复杂也是最重要的一种运算，矩阵乘法的概念是从实际需要中产生出来的，读者以后会看到这一点.

\noindent\textbf{定义 2.2.3}\quad 设有 $m\times k$ 矩阵 $A=(a_{ij})_{m\times k}$，以及 $k\times n$ 矩阵 $B=(b_{ij})_{k\times n}$. 定义 $A$ 和 $B$ 的乘积 $AB$ 是一个 $m\times n$ 矩阵且 $AB$ 的第 $(i,j)$ 元素
\[
c_{ij}=a_{i1}b_{1j}+a_{i2}b_{2j}+\cdots+a_{ik}b_{kj}.
\]

为了使读者看得更清楚，我们写出矩阵乘积的表达式：
\[
\left(\begin{array}{cccc}
a_{11} & a_{12} & \cdots & a_{1k}\\
a_{21} & a_{22} & \cdots & a_{2k}\\
\vdots & \vdots & & \vdots\\
a_{m1} & a_{m2} & \cdots & a_{mk}
\end{array}\right)
\left(\begin{array}{cccc}
b_{11} & b_{12} & \cdots & b_{1n}\\
b_{21} & b_{22} & \cdots & b_{2n}\\
\vdots & \vdots & & \vdots\\
b_{k1} & b_{k2} & \cdots & b_{kn}
\end{array}\right)
=
\left(\begin{array}{cccc}
\sum\limits_{r=1}^{k}a_{1r}b_{r1} & \sum\limits_{r=1}^{k}a_{1r}b_{r2} & \cdots & \sum\limits_{r=1}^{k}a_{1r}b_{rn}\\
\sum\limits_{r=1}^{k}a_{2r}b_{r1} & \sum\limits_{r=1}^{k}a_{2r}b_{r2} & \cdots & \sum\limits_{r=1}^{k}a_{2r}b_{rn}\\
\vdots & \vdots & & \vdots\\
\sum\limits_{r=1}^{k}a_{mr}b_{r1} & \sum\limits_{r=1}^{k}a_{mr}b_{r2} & \cdots & \sum\limits_{r=1}^{k}a_{mr}b_{rn}
\end{array}\right).
\]

为了掌握这个定义，我们需要注意以下两个要点：

第一，$A$ 和 $B$ 只有在 $A$ 的列数等于 $B$ 的行数时才可以相乘，这时的积写为 $AB$ (注意不能写为 $BA$). 得到的积矩阵 $AB$ 的行数等于 $A$ 的行数，列数等于 $B$ 的列数. 我们用下列式子来帮助记忆：
\[
\begin{array}{ccc}
A & B & \longrightarrow AB,\\
m\times k & k\times n & \longrightarrow m\times n,
\end{array}
\]
即 $AB$ 的行列数只需把 $A$ 与 $B$ 的行列数合并去掉中间的两个 $k$ 就可以了.

第二，$A$ 与 $B$ 的积 $AB$ 的第 $(i,j)$ 元素为
\[
c_{ij}=a_{i1}b_{1j}+a_{i2}b_{2j}+\cdots+a_{ik}b_{kj}.
\]
上式中只涉及 $A$ 的第 $i$ 行元素与 $B$ 的第 $j$ 列元素. 也就是说 $c_{ij}$ 等于将 $A$ 的第 $i$ 行元素与 $B$ 的第 $j$ 列元素分别相乘以后再求和的结果.

\noindent\textbf{例 2.2.1}\quad
\[
A=\left(\begin{array}{ccc}
1 & 0 & 1\\
2 & 1 & 0
\end{array}\right),\quad
B=\left(\begin{array}{cccc}
1 & 0 & 1 & 1\\
1 & 1 & 2 & -1\\
-1 & 0 & -1 & 0
\end{array}\right),
\]
求 $AB$.

\noindent\textbf{解}\quad $A$ 是 $2\times 3$ 矩阵，$B$ 是 $3\times 4$ 矩阵. 因此 $AB$ 是 $2\times 4$ 矩阵. 若令
\[
C=(c_{ij})=AB,
\]
则
\[
\begin{array}{rcl}
c_{11}&=&1\cdot 1+0\cdot 1+1\cdot(-1)=0,\\
c_{12}&=&1\cdot 0+0\cdot 1+1\cdot 0=0,\\
c_{13}&=&1\cdot 1+0\cdot 2+1\cdot(-1)=0,\\
c_{14}&=&1\cdot 1+0\cdot(-1)+1\cdot 0=1,\\
c_{21}&=&2\cdot 1+1\cdot 1+0\cdot(-1)=3,\\
c_{22}&=&2\cdot 0+1\cdot 1+0\cdot 0=1,\\
c_{23}&=&2\cdot 1+1\cdot 2+0\cdot(-1)=4,\\
c_{24}&=&2\cdot 1+1\cdot(-1)+0\cdot 0=1.
\end{array}
\]
所以
\[
AB=\left(\begin{array}{cccc}
0 & 0 & 0 & 1\\
3 & 1 & 4 & 1
\end{array}\right).
\]
在这个例子中，如果把 $A,B$ 的次序倒过来，则由于 $B$ 的列数等于 4, $A$ 的行数等于 2, $B$ 与 $A$ 不能相乘或者说它们的乘法无意义. 从这里可以看出，矩阵的乘法一般不满足交换律，即一般来说 $AB\ne BA$. 对于某些矩阵 $A,B$，即使 $AB$ 与 $BA$ 都有意义，它们也未必相等.

\noindent\textbf{例 2.2.2}\quad
\[
A=(1,0,4),\quad
B=\left(\begin{array}{c}
1\\
1\\
0
\end{array}\right),
\]
求 $AB$ 和 $BA$.

\noindent\textbf{解}\quad
\[
AB=(1,0,4)\left(\begin{array}{c}
1\\
1\\
0
\end{array}\right)=1\cdot 1+0\cdot 1+4\cdot 0=1,
\]
\[
BA=\left(\begin{array}{c}
1\\
1\\
0
\end{array}\right)(1,0,4)=
\left(\begin{array}{ccc}
1 & 0 & 4\\
1 & 0 & 4\\
0 & 0 & 0
\end{array}\right).
\]
从上面可以看出 $AB$ 是一个 $1\times 1$ 矩阵 (我们通常把 $1\times 1$ 矩阵看成是一个数，不必加括号)，而 $BA$ 是一个 $3\times 3$ 矩阵，显然 $AB\ne BA$.

\noindent\textbf{例 2.2.3}\quad
\[
A=\left(\begin{array}{cc}
0 & 0\\
0 & 1
\end{array}\right),\quad
B=\left(\begin{array}{cc}
0 & 1\\
0 & 0
\end{array}\right),
\]
求 $AB$ 和 $BA$.

\noindent\textbf{解}\quad
\[
AB=\left(\begin{array}{cc}
0 & 0\\
0 & 0
\end{array}\right),\quad
BA=\left(\begin{array}{cc}
0 & 1\\
0 & 0
\end{array}\right).
\]
在这个例子中，$A,B$ 都是二阶方阵，但 $AB$ 是零矩阵，$BA$ 不是零矩阵，因此 $AB\ne BA$. 由于矩阵乘法不满足交换律，因此读者在进行运算时千万要注意，不能把乘积的次序搞错.

矩阵的乘法比通常人们所熟悉的数的乘法要复杂得多. 读者也许会问：为什么要这样来定义矩阵乘法？当 $A=(a_{ij})_{m\times n}, B=(b_{ij})_{m\times n}$ 时，定义 $A$ 与 $B$ 之积为矩阵 $(a_{ij}\cdot b_{ij})_{m\times n}$ 不是更简单吗？我们说矩阵乘法的定义是出于实际需要. 将 $A$ 与 $B$ 的积定义成 $(a_{ij}\cdot b_{ij})_{m\times n}$ 在历史上也有过，但是这种乘法比起我们刚才定义的乘法来，用处要小得多. 矩阵乘法虽然比较复杂，但只要多练习，是不难掌握的. 下面我们举例来说明矩阵乘法的用途.

\noindent\textbf{例 2.2.4}\quad 设有 $n$ 个未知数 $m$ 个方程式的线性方程组：
\begin{equation}
\left\{\begin{array}{rcl}
a_{11}x_1+a_{12}x_2+\cdots+a_{1n}x_n&=&b_1,\\
a_{21}x_1+a_{22}x_2+\cdots+a_{2n}x_n&=&b_2,\\
\multicolumn{3}{c}{\cdots\cdots\cdots\cdots}\\
a_{m1}x_1+a_{m2}x_2+\cdots+a_{mn}x_n&=&b_m,
\end{array}\right.
\tag{2.2.1}
\end{equation}
令
\[
A=\left(\begin{array}{cccc}
a_{11} & a_{12} & \cdots & a_{1n}\\
a_{21} & a_{22} & \cdots & a_{2n}\\
\vdots & \vdots & & \vdots\\
a_{m1} & a_{m2} & \cdots & a_{mn}
\end{array}\right),
\]
\[
x=\left(\begin{array}{c}
x_1\\
x_2\\
\vdots\\
x_n
\end{array}\right),\quad
\beta=\left(\begin{array}{c}
b_1\\
b_2\\
\vdots\\
b_m
\end{array}\right).
\]
注意到 $A$ 是一个 $m\times n$ 矩阵，$x$ 是一个 $n\times 1$ 矩阵，因此 $Ax$ 是一个 $m\times 1$ 矩阵. 方程组 (2.2.1) 用矩阵乘法就可简写为
\begin{equation}
Ax=\beta.
\tag{2.2.2}
\end{equation}
这种写法不仅节省了篇幅，更重要的是可以用矩阵的方法来处理线性方程组.

\noindent\textbf{矩阵乘法运算规则}\quad 矩阵乘法运算适合下列规则：

(1) 结合律：$(AB)C=A(BC)$;

(2) 分配律：$A(B+C)=AB+AC,\ (A+B)C=AC+BC$;

(3) $c(AB)=(cA)B=A(cB)$，这里 $c$ 是一个数;

(4) 对任意的 $m\times n$ 矩阵 $A$, $I_mA=A=AI_n$，其中 $I_m$ 和 $I_n$ 分别是 $m$ 阶和 $n$ 阶单位阵，即单位阵在矩阵乘法运算中起的作用和数 $1$ 在数的乘法运算中起的作用类似.

\noindent\textbf{证明}\quad (1) 结合律.

首先我们注意到，由于 $A$ 与 $B$ 可相乘，$B$ 与 $C$ 也可相乘，因此可设 $A$ 是 $m\times n$ 阵，$B$ 是 $n\times p$ 阵，$C$ 是 $p\times q$ 阵. 于是，$AB$ 是一个 $m\times p$ 阵，$(AB)C$ 是一个 $m\times q$ 阵. 另一方面 $BC$ 是一个 $n\times q$ 阵，$A(BC)$ 是一个 $m\times q$ 阵. 因此只需验证两个 $m\times q$ 阵的任意第 $(i,j)$ 元素对应相等就可以了. 现设 $A=(a_{ij})_{m\times n}, B=(b_{ij})_{n\times p}, C=(c_{ij})_{p\times q}$. 矩阵 $(AB)C$ 的第 $(i,j)$ 元素可以看成是 $AB$ 的第 $i$ 行元素与 $C$ 的第 $j$ 列元素对应相乘之和，而 $AB$ 的第 $i$ 行的元素为
\[
\left(\sum_{r=1}^{n}a_{ir}b_{r1},\sum_{r=1}^{n}a_{ir}b_{r2},\cdots,\sum_{r=1}^{n}a_{ir}b_{rp}\right),
\]
$C$ 的第 $j$ 列为
\[
\left(\begin{array}{c}
c_{1j}\\
c_{2j}\\
\vdots\\
c_{pj}
\end{array}\right),
\]
因此，$(AB)C$ 的第 $(i,j)$ 元素为
\[
\left(\sum_{r=1}^{n}a_{ir}b_{r1}\right)c_{1j}+\left(\sum_{r=1}^{n}a_{ir}b_{r2}\right)c_{2j}+\cdots+\left(\sum_{r=1}^{n}a_{ir}b_{rp}\right)c_{pj}
\]
\[
=\sum_{k=1}^{p}\left(\sum_{r=1}^{n}a_{ir}b_{rk}\right)c_{kj}
=\sum_{k=1}^{p}\sum_{r=1}^{n}a_{ir}b_{rk}c_{kj}.
\]
另一方面，用同样办法可以求得 $A(BC)$ 的第 $(i,j)$ 元素为
\[
\sum_{r=1}^{n}a_{ir}\left(\sum_{k=1}^{p}b_{rk}c_{kj}\right)
=\sum_{r=1}^{n}\sum_{k=1}^{p}a_{ir}b_{rk}c_{kj}.
\]
由于
\[
\sum_{k=1}^{p}\sum_{r=1}^{n}a_{ir}b_{rk}c_{kj}
=
\sum_{r=1}^{n}\sum_{k=1}^{p}a_{ir}b_{rk}c_{kj},
\]
因此有
\[
(AB)C=A(BC).
\]

(2) 分配律.

设 $A=(a_{ij})_{m\times n}, B=(b_{ij})_{n\times p}, C=(c_{ij})_{n\times p}$，则 $B+C=(b_{ij}+c_{ij})_{n\times p}$. 因此 $A(B+C)$ 的第 $(i,j)$ 元素为
\[
a_{i1}(b_{1j}+c_{1j})+a_{i2}(b_{2j}+c_{2j})+\cdots+a_{in}(b_{nj}+c_{nj})
\]
\[
=(a_{i1}b_{1j}+a_{i2}b_{2j}+\cdots+a_{in}b_{nj})+(a_{i1}c_{1j}+a_{i2}c_{2j}+\cdots+a_{in}c_{nj}).
\]
显然这就是 $AB+AC$ 的第 $(i,j)$ 元素.

(3) 与 (4) 都很容易，请读者自己验证. $\square$

由矩阵乘法的结合律，我们今后将 $(AB)C$ 或 $A(BC)$ 直接写为 $ABC$，中间不再加括号. 利用结合律还可将若干个矩阵的乘积简写为 $A_1A_2\cdots A_m$，中间无需加任何括号.

我们还可以定义同一矩阵的乘方 (即幂). 记
\[
\begin{array}{rcl}
A^2&=&A\cdot A,\\
A^3&=&A\cdot A\cdot A,\\
\multicolumn{3}{c}{\cdots\cdots\cdots\cdots}\\
A^k&=&A\cdot A\cdots A\ (k\text{ 个 }A).
\end{array}
\]
要使 $A$ 的乘方有意义，$A$ 的行数必须等于列数，也就是说 $A$ 必须是方阵.

\noindent\textbf{方阵幂的运算规则}

(1) $A^rA^s=A^{r+s}$;

(2) $(A^r)^s=A^{rs}$.

由于矩阵的乘法不满足交换律，因此一般来说，若 $r>1$, $(AB)^r\ne A^rB^r$. 只有当 $AB=BA$ 时，才有 $(AB)^r=A^rB^r$. 虽然矩阵乘法不可交换，但对于某些特殊的矩阵，乘法仍是可交换的. 比如由矩阵乘法的运算规则 (4) 知，$I_n$ 与任一 $n$ 阶方阵乘法可交换. 事实上，若 $c$ 是某个数，则 $cI_n$ 与任一 $n$ 阶方阵乘法也可交换. 当 $AB=BA$ 时，不难验证``二项式定理''也成立：
\[
(A+B)^n=A^n+C_n^1A^{n-1}B+C_n^2A^{n-2}B^2+\cdots+C_n^{n-1}AB^{n-1}+B^n.
\]

矩阵乘法除了不可交换是它的一个特点外，还有一个特点是两个非零矩阵的乘积可能为零矩阵. 这一点由前面的例子我们已经看到. 由于这个缘故，对矩阵的乘法来说，消去律一般不成立，即若 $AB=AC$, $A\ne O$，我们不能推出 $B=C$. 何时可用消去律，我们将在后面讨论.

\subsection*{四、矩阵的转置}

矩阵的转置与行列式的转置定义是类似的.

\noindent\textbf{定义 2.2.4}\quad 设 $A=(a_{ij})$ 是 $m\times n$ 矩阵，定义 $A$ 的转置 $A'$ 为一个 $n\times m$ 矩阵，它的第 $k$ 行正好是矩阵 $A$ 的第 $k$ 列 $(k=1,2,\cdots,n)$; 它的第 $r$ 列是 $A$ 的第 $r$ 行 $(r=1,2,\cdots,m)$.

例如，设
\[
A=\left(\begin{array}{cccc}
1 & 2 & 1 & 0\\
-1 & 3 & 0 & 5\\
\sqrt{2} & 1 & -2 & 0
\end{array}\right),
\]
则
\[
A'=\left(\begin{array}{ccc}
1 & -1 & \sqrt{2}\\
2 & 3 & 1\\
1 & 0 & -2\\
0 & 5 & 0
\end{array}\right).
\]

\noindent\textbf{矩阵转置运算规则}

(1) $(A')'=A$;

(2) $(A+B)'=A'+B'$;

(3) $(cA)'=cA'$;

(4) $(AB)'=B'A'$.

\noindent\textbf{证明}\quad 上述 (1)--(3) 是显然的，现来证明 (4). 设 $A=(a_{ij})_{m\times n}, B=(b_{ij})_{n\times p}$，又设 $C=AB$，且 $C=(c_{ij})_{m\times p}$，则 $C$ 的第 $(i,j)$ 元素为
\[
c_{ij}=\sum_{r=1}^{n}a_{ir}b_{rj}.
\]
因此 $C'$ 的第 $(j,i)$ 元素为 $c_{ij}$. 再看 $B'A'$，它的第 $(j,i)$ 元素等于 $B'$ 的第 $j$ 行元素与 $A'$ 的第 $i$ 列元素对应相乘之和. 但 $B'$ 的第 $j$ 行元素等于 $B$ 的第 $j$ 列元素，$A'$ 的第 $i$ 列元素等于 $A$ 的第 $i$ 行元素. 它们对应元素相乘之和恰为
\[
a_{i1}b_{1j}+a_{i2}b_{2j}+\cdots+a_{in}b_{nj}=c_{ij}.
\]
另外，显然 $C'$ 与 $B'A'$ 的行、列数分别相等，因此 $C'=B'A'$. $\square$

一个矩阵经转置以后得到的矩阵一般来说与原矩阵不同. 如果一个方阵转置后仍与原矩阵相同，即 $A'=A$，则称这样的矩阵为对称阵. 例如，下列矩阵为对称阵：
\[
\left(\begin{array}{ccc}
1 & 1 & 0\\
1 & 3 & 5\\
0 & 5 & -1
\end{array}\right).
\]
若一个方阵经转置后等于原矩阵的负矩阵，即 $A'=-A$，就称它是一个反对称阵. 例如，下列矩阵为反对称阵：
\[
\left(\begin{array}{ccc}
0 & 2 & 1\\
-2 & 0 & -5\\
-1 & 5 & 0
\end{array}\right).
\]
对称阵特别是实对称阵我们今后还要进行仔细研究. 读者不难看出，对称阵的元素以主对角线为对称线，即 $a_{ij}=a_{ji}$; 反对称阵主对角线上的元素皆为零，且 $a_{ij}=-a_{ji}$.

\subsection*{五、矩阵的共轭}

复矩阵还有一种运算，称为共轭运算. 设 $z$ 是一个复数，$z=a+bi$，我们用 $\bar z$ 表示 $z$ 的共轭复数 $a-bi$.

\noindent\textbf{定义 2.2.5}\quad 设 $A=(a_{ij})_{m\times n}$ 是一个复矩阵，则 $A$ 的共轭矩阵 $\overline{A}$ 是一个 $m\times n$ 复矩阵，且
\[
\overline{A}=(\overline{a_{ij}})_{m\times n}.
\]
这就是说 $A$ 的共轭矩阵的每个第 $(i,j)$ 元素是 $A$ 的第 $(i,j)$ 元素的共轭复数. 例如
\[
A=\left(\begin{array}{ccc}
1+i & 0 & 1+\sqrt{2}i\\
2i & -1 & -4i\\
-4-i & \sqrt{3}i & i
\end{array}\right),
\]
则
\[
\overline{A}=\left(\begin{array}{ccc}
1-i & 0 & 1-\sqrt{2}i\\
-2i & -1 & 4i\\
-4+i & -\sqrt{3}i & -i
\end{array}\right).
\]

\noindent\textbf{矩阵共轭运算规则}

(1) $\overline{A+B}=\overline{A}+\overline{B}$;

(2) $\overline{cA}=\bar c\,\overline{A}$;

(3) $\overline{AB}=\overline{A}\,\overline{B}$;

(4) $\overline{(A')}=(\overline{A})'$.

这些规则很容易验证，请读者自己完成.

\begin{center}
\textbf{习\quad 题\quad 2.2}
\end{center}

1. 计算：
\[
(1)\ \left(\begin{array}{cc}
1 & 2\\
-3 & 0
\end{array}\right)+
\left(\begin{array}{cc}
2 & -2\\
0 & 1
\end{array}\right);
\qquad
(2)\ 3\left(\begin{array}{cc}
1 & \sqrt{2}\\
0 & -4
\end{array}\right);
\]
\[
(3)\ 2\left(\begin{array}{ccc}
1 & 2 & 3\\
4 & -1 & \sqrt{2}
\end{array}\right)
-3\left(\begin{array}{ccc}
\frac{1}{3} & 1 & 0\\
0 & -1 & -\sqrt{2}
\end{array}\right);
\]
\[
(4)\ \left(\begin{array}{cccc}
1 & 2 & 3 & 4\\
0 & 1 & -1 & 0\\
3 & 4 & 1 & 5
\end{array}\right)
-\frac{1}{2}\left(\begin{array}{cccc}
2 & 4 & 0 & 1\\
0 & 2 & 3 & 1\\
4 & 2 & 8 & 0
\end{array}\right).
\]

2. 计算：
\[
(1)\ \left(\begin{array}{ccc}
1 & 2 & 0\\
1 & -1 & 1
\end{array}\right)
\left(\begin{array}{cc}
1 & 3\\
0 & 1\\
1 & -1
\end{array}\right);
\qquad
(2)\ \left(\begin{array}{cc}
2 & 0\\
0 & 1
\end{array}\right)
\left(\begin{array}{cc}
-1 & 1\\
-2 & 0
\end{array}\right);
\]
\[
(3)\ \left(\begin{array}{cccc}
0 & 1 & -1 & 3\\
-1 & 2 & 1 & 0
\end{array}\right)
\left(\begin{array}{cc}
1 & 1\\
-1 & 4\\
3 & 0\\
1 & 2
\end{array}\right);
\qquad
(4)\ \left(\begin{array}{ccc}
2 & 1 & -2\\
1 & 0 & 4\\
-3 & 1 & 0\\
0 & 1 & 1
\end{array}\right)
\left(\begin{array}{ccc}
3 & 1 & 0\\
0 & 0 & 1\\
-1 & 2 & 0
\end{array}\right).
\]

3. 设
\[
A=\left(\begin{array}{ccc}
1 & 0 & 3\\
2 & -1 & 0
\end{array}\right),\quad
B=\left(\begin{array}{cc}
1 & -1\\
2 & 3\\
4 & 0
\end{array}\right),
\]
试求 $AB$ 与 $BA$.

4. 计算下列矩阵的 $k$ 次幂，其中 $k$ 为正整数：
\[
(1)\ A=\left(\begin{array}{ccc}
a & 0 & 0\\
0 & b & 0\\
0 & 0 & c
\end{array}\right);
\qquad
(2)\ A=\left(\begin{array}{cc}
\cos\theta & -\sin\theta\\
\sin\theta & \cos\theta
\end{array}\right);
\]
\[
(3)\ A=\left(\begin{array}{ccc}
a & 1 & 0\\
0 & a & 1\\
0 & 0 & a
\end{array}\right);
\qquad
(4)\ A=\left(\begin{array}{ccc}
1 & 2 & 4\\
2 & 4 & 8\\
3 & 6 & 12
\end{array}\right).
\]

5. 设 $A=(a_{ij})$ 是 $n$ 阶对称阵，$x=(x_1,x_2,\cdots,x_n)'$ 是 $n$ 维列向量，试求 $x'Ax$.

6. 求证：上 (下) 三角阵的和、差、数乘及乘积仍是上 (下) 三角阵，并且所得上 (下) 三角阵的主对角元是原上 (下) 三角阵的对应主对角元的和、差、数乘及乘积.

7. 若 $A$ 是 $n$ 阶方阵且 $A^n=O$，求证：
\[
(I_n-A)(I_n+A+A^2+\cdots+A^{n-1})=I_n.
\]

8. 设 $A$ 是实对称阵，若 $A^2=O$，求证：$A=O$.

9. 证明：任一 $n$ 阶方阵均可表示为一个对称阵和一个反对称阵之和.

10. 设 $A$ 是 $n$ 阶复矩阵，若 $\overline{A'}=A$ (即矩阵的共轭转置等于自身)，则称 $A$ 是一个 Hermite (厄米特) 矩阵. 若 $\overline{A'}=-A$，称 $A$ 是斜 Hermite 矩阵. 求证：任一 $n$ 阶复矩阵均可表示为一个 Hermite 矩阵与一个斜 Hermite 矩阵之和.

11. 设 $A$ 是 $n$ 阶复矩阵，求证：$|\overline{A}|=\overline{|A|}$.

12. 设 $A,B$ 为 $n$ 阶方阵，求证：

(1) 若 $A,B$ 为对称阵，则 $AB$ 为对称阵的充分必要条件是 $AB=BA$, $AB$ 为反对称阵的充分必要条件是 $AB=-BA$;

(2) 若 $A$ 为对称阵，$B$ 为反对称阵，则 $AB$ 为反对称阵的充分必要条件是 $AB=BA$, $AB$ 为对称阵的充分必要条件是 $AB=-BA$.

13. 求与矩阵 $A$ 乘法可交换的所有矩阵：
\[
(1)\ A=\left(\begin{array}{cc}
1 & 1\\
0 & 1
\end{array}\right);
\qquad
(2)\ A=\left(\begin{array}{cc}
0 & 1\\
1 & 0
\end{array}\right);
\]
\[
(3)\ A=\left(\begin{array}{ccc}
1 & 0 & 0\\
0 & 1 & 0\\
3 & 1 & 1
\end{array}\right);
\qquad
(4)\ A=\left(\begin{array}{ccc}
0 & 1 & 0\\
0 & 0 & 1\\
1 & 1 & 0
\end{array}\right).
\]

14. 求证：

(1) 与所有 $n$ 阶对角阵乘法可交换的矩阵也必是 $n$ 阶对角阵;

(2) 与所有 $n$ 阶矩阵乘法可交换的矩阵是形如 $cI_n$ 的对角阵 ($cI_n$ 称为纯量阵或数量阵).

\section*{§2.3\quad 方阵的逆阵}

我们已经介绍了矩阵的加减法、数乘、乘法等运算，读者会问：矩阵有除法吗？我们知道，在数的运算中，除法是乘法的逆运算，它可以通过求倒数来实现. 即若求数 $a$ 除以 $b\ (b\ne 0)$ 的商，只需求出 $b^{-1}=\frac{1}{b}$，则 $\frac{a}{b}=ab^{-1}$. 对矩阵我们也可以这样做，先定义矩阵的逆阵，然后将矩阵的除法归结为一个矩阵和另外一个矩阵的逆阵之积. 但是现在有一个问题：矩阵的乘法一般是不可交换的，因此对矩阵逆的定义要有更严格的要求.

\noindent\textbf{定义 2.3.1}\quad 设 $A$ 是 $n$ 阶方阵，若存在一个 $n$ 阶方阵 $B$，使得
\[
AB=BA=I_n,
\]
则称 $B$ 是 $A$ 的逆阵，记为 $B=A^{-1}$. 凡有逆阵的矩阵称为可逆阵或非奇异阵 (简称非异阵)，否则称为奇异阵.

矩阵的求逆运算有它自己的特点，我们必须注意：

(1) 根据上述定义，只有对方阵才有逆阵的定义，长方阵没有逆阵.

(2) 并非任一非零方阵都有逆阵. 比如，矩阵
\[
A=\left(\begin{array}{cc}
1 & 1\\
0 & 0
\end{array}\right)
\]
就没有逆阵. 因为对任一 $B=(b_{ij})_{2\times 2}$，有
\[
AB=\left(\begin{array}{cc}
1 & 1\\
0 & 0
\end{array}\right)
\left(\begin{array}{cc}
b_{11} & b_{12}\\
b_{21} & b_{22}
\end{array}\right)
=
\left(\begin{array}{cc}
b_{11}+b_{21} & b_{12}+b_{22}\\
0 & 0
\end{array}\right).
\]
$AB$ 不可能是单位阵.

\noindent\textbf{注}\quad 从这里我们可以发现，若某个矩阵有一行元素全等于零，则这个矩阵一定不是可逆阵. 同理不难证明，若某个矩阵的一列元素全等于零，则该矩阵也必不是可逆阵.

(3) 由于矩阵的乘法一般不满足交换律，因此一般来说，$AB^{-1}\ne B^{-1}A$，即右除不一定等于左除.

一个 $n$ 阶方阵 $A$ 若有逆阵，则逆阵唯一吗？设 $B,C$ 是 $n$ 阶方阵，且都是 $A$ 的逆阵，即
\[
AB=BA=I_n,\quad AC=CA=I_n,
\]
则
\[
B=BI_n=B(AC)=(BA)C=I_nC=C.
\]
因此只要矩阵可逆，其逆阵必唯一.

\noindent\textbf{矩阵求逆运算规则}

(1) 若 $A$ 是非异阵，则 $(A^{-1})^{-1}=A$;

(2) 若 $A,B$ 都是 $n$ 阶非异阵，则 $AB$ 也是 $n$ 阶非异阵且 $(AB)^{-1}=B^{-1}A^{-1}$;

(3) 若 $A$ 是非异阵，$c$ 是非零数，则 $cA$ 也是非异阵且 $(cA)^{-1}=c^{-1}A^{-1}$;

(4) 若 $A$ 是非异阵，则 $A$ 的转置 $A'$ 也是非异阵且 $(A')^{-1}=(A^{-1})'$.

\noindent\textbf{证明}\quad (1) 因为 $(A^{-1})A=A(A^{-1})=I_n$，故 $A$ 是 $A^{-1}$ 的逆阵.

(2) 我们有
\[
(AB)(B^{-1}A^{-1})=A(BB^{-1})A^{-1}=AI_nA^{-1}=AA^{-1}=I_n.
\]
同理 $(B^{-1}A^{-1})(AB)=I_n$.

(3) 显然.

(4) 由 $AA^{-1}=I_n$，两边转置并注意到 $I_n'=I_n$，得
\[
(AA^{-1})'=I_n.
\]
由转置的性质知
\[
(AA^{-1})'=(A^{-1})'A'.
\]
因此
\[
(A^{-1})'A'=I_n.
\]
同理可得
\[
A'(A^{-1})'=I_n.
\]

需要提醒读者注意的是规则 (2)：$(AB)^{-1}=B^{-1}A^{-1}$. 一般来说，$(AB)^{-1}\ne A^{-1}B^{-1}$，只有当 $AB=BA$ 时，才有 $(AB)^{-1}=A^{-1}B^{-1}$.

用数学归纳法容易证明：
\[
(A_1A_2\cdots A_k)^{-1}=A_k^{-1}\cdots A_2^{-1}A_1^{-1},
\]
其中每个 $A_i$ 都是可逆阵.

如何求逆阵？我们介绍一个方法.

设 $A$ 是 $n$ 阶方阵，这个方阵决定了一个 $n$ 阶行列式，记为 $|A|$ 或 $\det A$. 若 $A=(a_{ij})_{n\times n}$，则
\[
|A|=\left|\begin{array}{cccc}
a_{11} & a_{12} & \cdots & a_{1n}\\
a_{21} & a_{22} & \cdots & a_{2n}\\
\vdots & \vdots & & \vdots\\
a_{n1} & a_{n2} & \cdots & a_{nn}
\end{array}\right|.
\]

\noindent\textbf{定义 2.3.2}\quad 设 $A$ 是 $n$ 阶方阵，$A_{ij}$ 是行列式 $|A|$ 中第 $(i,j)$ 元素 $a_{ij}$ 的代数余子式，则称下列方阵为 $A$ 的伴随阵：
\[
\left(\begin{array}{cccc}
A_{11} & A_{21} & \cdots & A_{n1}\\
A_{12} & A_{22} & \cdots & A_{n2}\\
\vdots & \vdots & & \vdots\\
A_{1n} & A_{2n} & \cdots & A_{nn}
\end{array}\right).
\]
$A$ 的伴随阵通常记为 $A^*$.

方阵与其伴随阵之间有如下的关系：

\noindent\textbf{引理 2.3.1}\quad 设 $A$ 为 $n$ 阶方阵，$A^*$ 为 $A$ 的伴随阵，则
\[
AA^*=A^*A=|A|\cdot I_n.
\]

\noindent\textbf{证明}\quad
\[
AA^*=
\left(\begin{array}{cccc}
a_{11} & a_{12} & \cdots & a_{1n}\\
a_{21} & a_{22} & \cdots & a_{2n}\\
\vdots & \vdots & & \vdots\\
a_{n1} & a_{n2} & \cdots & a_{nn}
\end{array}\right)
\left(\begin{array}{cccc}
A_{11} & A_{21} & \cdots & A_{n1}\\
A_{12} & A_{22} & \cdots & A_{n2}\\
\vdots & \vdots & & \vdots\\
A_{1n} & A_{2n} & \cdots & A_{nn}
\end{array}\right).
\]
上式中两个矩阵乘积的第 $(i,j)$ 元素为
\[
a_{i1}A_{j1}+a_{i2}A_{j2}+\cdots+a_{in}A_{jn}.
\]
由定理 1.4.2 可知，当 $i=j$ 时上式为 $|A|$，当 $i\ne j$ 时上式等于零. 因此
\[
AA^*=
\left(\begin{array}{cccc}
|A| & 0 & \cdots & 0\\
0 & |A| & \cdots & 0\\
\vdots & \vdots & & \vdots\\
0 & 0 & \cdots & |A|
\end{array}\right)
=|A|\cdot I_n.
\]
由定理 1.4.1 同理可证
\[
A^*A=|A|\cdot I_n.
\]

\noindent\textbf{定理 2.3.1}\quad 若 $|A|\ne 0$，则 $A$ 是一个非异阵，且
\[
A^{-1}=\frac{1}{|A|}A^*.
\]

\noindent\textbf{证明}\quad 注意到 $|A|\ne 0$，由引理 2.3.1 可得
\[
A\left(\frac{1}{|A|}A^*\right)=\left(\frac{1}{|A|}A^*\right)A=I_n,
\]
因此 $A$ 是一个非异阵，且 $A^{-1}=\frac{1}{|A|}A^*$. $\square$

\noindent\textbf{注}\quad 我们在以后将证明若 $A$ 是非异阵，则 $|A|\ne 0$. 因此上述定理提供了计算逆阵的一般方法. 但是用这个定理计算逆阵必须计算所有的 $A_{ij}$，计算量一般是相当大的，我们将在 §2.5 中介绍一种比较简单的计算方法.

利用逆阵来解线性方程组将显得特别简单. 由 §2.2 知道一个线性方程组
\[
\left\{\begin{array}{rcl}
a_{11}x_1+a_{12}x_2+\cdots+a_{1n}x_n&=&b_1,\\
a_{21}x_1+a_{22}x_2+\cdots+a_{2n}x_n&=&b_2,\\
\multicolumn{3}{c}{\cdots\cdots\cdots\cdots}\\
a_{n1}x_1+a_{n2}x_2+\cdots+a_{nn}x_n&=&b_n
\end{array}\right.
\]
可写成矩阵形式
\[
Ax=\beta,
\]
其中 $A=(a_{ij})$. 若 $|A|\ne 0$，则 $A^{-1}$ 必存在，因此
\[
A^{-1}(Ax)=A^{-1}\beta,
\]
即
\[
x=A^{-1}\beta.
\]
将上式中的矩阵写出来就是：
\[
\left(\begin{array}{c}
x_1\\
x_2\\
\vdots\\
x_n
\end{array}\right)
=\frac{1}{|A|}
\left(\begin{array}{cccc}
A_{11} & A_{21} & \cdots & A_{n1}\\
A_{12} & A_{22} & \cdots & A_{n2}\\
\vdots & \vdots & & \vdots\\
A_{1n} & A_{2n} & \cdots & A_{nn}
\end{array}\right)
\left(\begin{array}{c}
b_1\\
b_2\\
\vdots\\
b_n
\end{array}\right).
\]
于是
\[
x_1=\frac{1}{|A|}(b_1A_{11}+b_2A_{21}+\cdots+b_nA_{n1})=\frac{|A_1|}{|A|},
\]
其中
\[
|A_1|=\left|\begin{array}{cccc}
b_1 & a_{12} & \cdots & a_{1n}\\
b_2 & a_{22} & \cdots & a_{2n}\\
\vdots & \vdots & & \vdots\\
b_n & a_{n2} & \cdots & a_{nn}
\end{array}\right|.
\]
同理可得其余的 $x_i$. 显然这就是 Cramer 法则.

\begin{center}
\textbf{习\quad 题\quad 2.3}
\end{center}

1. 求下列矩阵的逆阵：
\[
(1)\ A=\left(\begin{array}{ccc}
1 & 1 & -1\\
1 & 2 & -3\\
0 & 1 & 1
\end{array}\right);
\qquad
(2)\ A=\left(\begin{array}{cccc}
a_1 & 0 & \cdots & 0\\
0 & a_2 & \cdots & 0\\
\vdots & \vdots & & \vdots\\
0 & 0 & \cdots & a_n
\end{array}\right)\ (a_i\ne 0).
\]

2. 用求逆阵的方法解线性方程组：
\[
\left\{\begin{array}{rcl}
x_1-x_2+3x_3&=&8,\\
2x_1-x_2+4x_3&=&11,\\
-x_1+2x_2-4x_3&=&-11.
\end{array}\right.
\]

3. 求证：有一行元素或一列元素全为零的 $n$ 阶方阵必是奇异阵.

4. 设 $A$ 是非异阵，求证：对任一正整数 $k$，有 $(A^k)^{-1}=(A^{-1})^k$.

5. 设 $A$ 是非异阵，证明乘法消去律成立，即从 $AB=AC$，可推出 $B=C$; 从 $BA=CA$ 也可推出 $B=C$.

6. 若 $n$ 阶方阵 $A$ 满足 $A^2=I_n$，则称为对合阵. 设 $A$ 是对合阵且 $I_n+A$ 是非异阵，求证：$A=I_n$.

7. 设 $A$ 是非零实矩阵且 $A^*=A'$，求证：$A$ 是非异阵.

8. 设 $A,B$ 及 $A+B$ 都是非异阵，求证：$A^{-1}+B^{-1}$ 也是非异阵.

9. 若 $n$ 阶方阵 $A$ 满足 $A^m=O$，其中 $m$ 为正整数，则称为幂零阵. 设 $A$ 是幂零阵，求证：$I_n-A$ 是非异阵.

10. 若 $n$ 阶方阵 $A$ 满足 $A^2=A$，则称为幂等阵. 设 $A$ 是幂等阵，求证：$A+I_n$ 是非异阵.

11. 设 $A$ 适合条件 $A^2-A-3I_n=O$，求证：$A-2I_n$ 是非异阵.

12. 设 $n$ 阶方阵 $A,B$ 满足 $A+B=AB$，求证：$I_n-A$ 是可逆阵且 $AB=BA$.

\section*{§2.4\quad 矩阵的初等变换与初等矩阵}

\subsection*{一、Gauss 消去法与矩阵的初等变换}

用 Gauss 消去法来解线性方程组是一种简便常用的方法，特别当未知数个数
不太多时，人们乐意采用这种方法在计算机上解线性方程组. 我们举例来说明这种方法.

\noindent\textbf{例 2.4.1}\quad 用 Gauss 消去法求解下列线性方程组：
\[
\left\{\begin{array}{rcl}
x_2+x_3&=&2,\\
2x_1+3x_2+2x_3&=&5,\\
3x_1+x_2-x_3&=&-1.
\end{array}\right.
\]

\noindent\textbf{解}\quad 将上述方程组中的第一式与第二式对调得
\[
\left\{\begin{array}{rcl}
2x_1+3x_2+2x_3&=&5,\\
x_2+x_3&=&2,\\
3x_1+x_2-x_3&=&-1.
\end{array}\right.
\]
将得到的方程组的第一式乘以 $-\frac{3}{2}$ 加到第三式上：
\[
\left\{\begin{array}{rcl}
2x_1+3x_2+2x_3&=&5,\\
x_2+x_3&=&2,\\
-\frac{7}{2}x_2-4x_3&=&-\frac{17}{2}.
\end{array}\right.
\]
将上面的第二式乘以 $\frac{7}{2}$ 加到第三式上：
\[
\left\{\begin{array}{rcl}
2x_1+3x_2+2x_3&=&5,\\
x_2+x_3&=&2,\\
-\frac{1}{2}x_3&=&-\frac{3}{2}.
\end{array}\right.
\]
将得到的第三式两边乘以 $-2$:
\[
\left\{\begin{array}{rcl}
2x_1+3x_2+2x_3&=&5,\\
x_2+x_3&=&2,\\
x_3&=&3.
\end{array}\right.
\]
将上面的最后一式分别乘以 $-1,-2$ 加到第二式及第一式上：
\[
\left\{\begin{array}{rcl}
2x_1+3x_2&=&-1,\\
x_2&=&-1,\\
x_3&=&3.
\end{array}\right.
\]
最后将上面的第二式乘以 $-3$ 加到第一式上，两边再乘以 $\frac{1}{2}$ 就得到方程组的解：
\[
\left\{\begin{array}{rcl}
x_1&=&1,\\
x_2&=&-1,\\
x_3&=&3.
\end{array}\right.
\]

上述求解过程可以用矩阵的变换来代替，将原方程组的系数及常数列排成一个矩阵，称为系数矩阵的增广矩阵，用 $\widetilde{A}$ 表示：
\[
\widetilde{A}=
\left(\begin{array}{cccc}
0 & 1 & 1 & 2\\
2 & 3 & 2 & 5\\
3 & 1 & -1 & -1
\end{array}\right).
\]
现在将求解过程用矩阵来描述如下.

将 $\widetilde{A}$ 的第一行与第二行对换得
\[
\left(\begin{array}{cccc}
2 & 3 & 2 & 5\\
0 & 1 & 1 & 2\\
3 & 1 & -1 & -1
\end{array}\right).
\]
将上面矩阵的第一行乘以 $-\frac{3}{2}$ 加到第三行上：
\[
\left(\begin{array}{cccc}
2 & 3 & 2 & 5\\
0 & 1 & 1 & 2\\
0 & -\frac{7}{2} & -4 & -\frac{17}{2}
\end{array}\right).
\]
将第二行乘以 $\frac{7}{2}$ 加到第三行上：
\[
\left(\begin{array}{cccc}
2 & 3 & 2 & 5\\
0 & 1 & 1 & 2\\
0 & 0 & -\frac{1}{2} & -\frac{3}{2}
\end{array}\right).
\]
将第三行乘以 $-2$ 得
\[
\left(\begin{array}{cccc}
2 & 3 & 2 & 5\\
0 & 1 & 1 & 2\\
0 & 0 & 1 & 3
\end{array}\right).
\]
将第三行乘以 $-1$ 加到第二行上，将第三行乘以 $-2$ 加到第一行上得
\[
\left(\begin{array}{cccc}
2 & 3 & 0 & -1\\
0 & 1 & 0 & -1\\
0 & 0 & 1 & 3
\end{array}\right).
\]
将第二行乘以 $-3$ 加到第一行上，最后将第一行乘以 $\frac{1}{2}$:
\[
\left(\begin{array}{cccc}
1 & 0 & 0 & 1\\
0 & 1 & 0 & -1\\
0 & 0 & 1 & 3
\end{array}\right).
\]

从上面的分析可以看出，用 Gauss 消去法解线性方程组的过程可以归结为对矩阵的变换，这不仅简化了解方程组的过程，更重要的是为彻底弄清线性方程组解的理论提供了工具. 上面例子适用于求解一般的线性方程组，我们把这种矩阵形式的 Gauss 消去法归结如下：

第一步：将线性方程组写成标准形式并写出系数矩阵的增广矩阵 $\widetilde{A}$.

第二步：将 $\widetilde{A}$ 中某一行调到第一行，使第一行第一列的元素不为零.

第三步：将得到的矩阵的第一行乘以某个数加到第二行上，消去第二行第一列的元素. 重复这一方法，直到消去第一列除第一行以外的所有元素.

第四步：重复上述步骤，使第二行第二列的元素不为零并消去第二列上其余元素. 不断用这个方法，将系数矩阵变成对角阵.

第五步：在每一行乘以某个非零数使系数矩阵变为单位阵，从而写出线性方程组的解.

在上述步骤中，我们对矩阵施行了以下 3 种变换：

(1) 两行对换;

(2) 以某一非零数乘以某一行;

(3) 以某一数乘以某一行后加到另一行上去.

这 3 种变换并不改变线性方程组的解. 也就是说，对应的新方程组与原方程组总是同解的.

\noindent\textbf{定义 2.4.1}\quad 下列 3 种矩阵变换分别称为矩阵的第一类、第二类、第三类初等行 (列) 变换：

(1) 对调矩阵中某两行 (列) 的位置;

(2) 用一非零常数 $c$ 乘以矩阵的某一行 (列);

(3) 将矩阵的某一行 (列) 乘以常数 $c$ 后加到另一行 (列) 上去.

上述 3 种变换统称为矩阵的初等变换.

显而易见，我们在上面对方程组的增广矩阵实施了矩阵的初等行变换.

\noindent\textbf{定义 2.4.2}\quad 如果一个矩阵 $A$ 经过有限次初等变换后变成 $B$，则称 $A$ 与 $B$ 是等价的，或 $A$ 与 $B$ 相抵，记为 $A\sim B$.

用初等变换可以将一个矩阵化简到什么程度？下面的定理回答了这个问题.

\noindent\textbf{定理 2.4.1}\quad 任一 $m\times n$ 矩阵 $A=(a_{ij})_{m\times n}$ 必相抵于下列 $m\times n$ 矩阵：
\begin{equation}
B=\left(\begin{array}{cccccc}
1 & \cdots & 0 & 0 & \cdots & 0\\
\vdots & & \vdots & \vdots & & \vdots\\
0 & \cdots & 1 & 0 & \cdots & 0\\
0 & \cdots & 0 & 0 & \cdots & 0\\
\vdots & & \vdots & \vdots & & \vdots\\
0 & \cdots & 0 & 0 & \cdots & 0
\end{array}\right),
\tag{2.4.1}
\end{equation}
其中 $B$ 的前 $r$ 行及前 $r$ 列交点处有 $r$ 个 1, 其余元素皆为零. 换言之，任一 $m\times n$ 矩阵均与一个主对角线上元素等于 1 或 0 而其余元素均为 0 的 $m\times n$ 矩阵相抵.

\noindent\textbf{证明}\quad 若 $A=O$，则结论显然成立. 现设 $A\ne O$，即 $A$ 至少有一个元素 $a_{ij}\ne 0$. 如果 $a_{ij}$ 不在第 $(1,1)$ 位置，那么可将它所在的行与第一行对换，再将它所在的列与第一列对换就可将 $a_{ij}$ 调至第 $(1,1)$ 位置. 因此我们不妨设 $a_{11}\ne 0$. 接下去将第一行依次乘以 $-a_{11}^{-1}a_{i1}$ 加到第 $i$ 行上去 $(i=2,3,\cdots,m)$，于是第一列元素除 $a_{11}$ 外都变成了零. 再将第一列元素乘以 $-a_{11}^{-1}a_{1j}$ 后加到第 $j$ 列上去 $(j=2,3,\cdots,n)$，则第一行元素除了 $a_{11}$ 外都变成零. 再用 $a_{11}^{-1}$ 乘以第一行，就得到第 $(1,1)$ 元素等于 1 而第一行及第一列其他元素都是零的矩阵，其形状如下：
\[
\left(\begin{array}{cccc}
1 & 0 & \cdots & 0\\
0 & b_{22} & \cdots & b_{2n}\\
\vdots & \vdots & & \vdots\\
0 & b_{m2} & \cdots & b_{mn}
\end{array}\right).
\]
再对第二行第二列采用与上面相同的步骤使第 $(2,2)$ 位置的元素不等于 0，并用同样办法消去除第 $(2,2)$ 元素外第二行及第二列的所有元素. 显然在进行上述过程中第一列及第一行的元素保持不变. 这样不断做下去，直到变成 (2.4.1) 式的形状为止. $\square$

\noindent\textbf{注}\quad 在这个定理中，对给定的矩阵 $A$，无论进行怎样的初等变换，最后得到的对角形矩阵中的 $r$ 总是不变的，这一重要事实将在第三章予以证明. 矩阵 (2.4.1) 称为矩阵 $A$ 的相抵标准型.

\noindent\textbf{例 2.4.2}\quad 用初等变换将下列矩阵化为相抵标准型：
\[
\left(\begin{array}{cccc}
2 & 0 & -1 & 3\\
1 & 2 & -2 & 4\\
0 & 1 & 3 & -1
\end{array}\right).
\]

\noindent\textbf{解}\quad
\[
\left(\begin{array}{cccc}
2 & 0 & -1 & 3\\
1 & 2 & -2 & 4\\
0 & 1 & 3 & -1
\end{array}\right)
\longrightarrow
\left(\begin{array}{cccc}
1 & 2 & -2 & 4\\
2 & 0 & -1 & 3\\
0 & 1 & 3 & -1
\end{array}\right)
\longrightarrow
\left(\begin{array}{cccc}
1 & 2 & -2 & 4\\
0 & -4 & 3 & -5\\
0 & 1 & 3 & -1
\end{array}\right).
\]
化到这一步接下去可将第一列分别乘以 $-2$、乘以 2、乘以 $-4$ 后加到第二列、第三列、第四列上去，使第一行中除第 $(1,1)$ 元素外其余都等于零. 但这时由于第一列除第一个元素外其余都为零，因此在整个过程中第二列、第三列、第四列的元素除处在第一行的元素外都不变，因此我们不必再写出具体过程而直接写出结果：
\[
\left(\begin{array}{cccc}
1 & 2 & -2 & 4\\
0 & -4 & 3 & -5\\
0 & 1 & 3 & -1
\end{array}\right)
\longrightarrow
\left(\begin{array}{cccc}
1 & 0 & 0 & 0\\
0 & -4 & 3 & -5\\
0 & 1 & 3 & -1
\end{array}\right).
\]
接下去再继续进行初等变换：
\[
\left(\begin{array}{cccc}
1 & 0 & 0 & 0\\
0 & -4 & 3 & -5\\
0 & 1 & 3 & -1
\end{array}\right)
\longrightarrow
\left(\begin{array}{cccc}
1 & 0 & 0 & 0\\
0 & 1 & 3 & -1\\
0 & -4 & 3 & -5
\end{array}\right).
\]
这一步是为了避免分数运算.
\[
\left(\begin{array}{cccc}
1 & 0 & 0 & 0\\
0 & 1 & 3 & -1\\
0 & -4 & 3 & -5
\end{array}\right)
\longrightarrow
\left(\begin{array}{cccc}
1 & 0 & 0 & 0\\
0 & 1 & 3 & -1\\
0 & 0 & 15 & -9
\end{array}\right)
\longrightarrow
\left(\begin{array}{cccc}
1 & 0 & 0 & 0\\
0 & 1 & 0 & 0\\
0 & 0 & 15 & -9
\end{array}\right)
\longrightarrow
\left(\begin{array}{cccc}
1 & 0 & 0 & 0\\
0 & 1 & 0 & 0\\
0 & 0 & 15 & 0
\end{array}\right)
\longrightarrow
\left(\begin{array}{cccc}
1 & 0 & 0 & 0\\
0 & 1 & 0 & 0\\
0 & 0 & 1 & 0
\end{array}\right).
\]
在进行初等变换时，有些步骤可并在一起以节省篇幅.

我们在进行初等变换的过程中，通常需要同时用行变换和列变换，如果将初等变换仅限制为行变换，定理 2.4.1 还成立吗？

考虑下面的 $1\times 3$ 矩阵:
\[
(0,1,2).
\]
显然，我们仅用初等行变换无论如何也不能将它化为标准型. 但是我们可以用初等行变换将矩阵化为所谓的阶梯形 (上阶梯形).

\noindent\textbf{定义 2.4.3}\quad 设 $A=(a_{ij})_{m\times n}$ 为 $m\times n$ 矩阵，对任意的 $1\le i\le m$，定义 $k_i$ 如下：若 $A$ 的第 $i$ 行元素全为零，则 $k_i=+\infty$；若 $A$ 的第 $i$ 行元素不全为零，则 $k_i$ 是第 $i$ 行所有非零元素列指标的最小值. 即若 $k_i<+\infty$，则 $a_{ik_i}$ 是 $A$ 的第 $i$ 行中从左至右第一个非零元素，$a_{ik_i}$ 称为第 $i$ 行的阶梯点.

若存在 $0\le r\le m$，使得 $k_1<k_2<\cdots<k_r,\ k_{r+1}=\cdots=k_m=+\infty$，则称这样的矩阵 $A$ 为阶梯形矩阵. 简单地说，若矩阵阶梯点的列指标随着行数严格递增，或者从图形上看非零元素全体构成一个阶梯，我们就称之为阶梯形矩阵.

下面几个矩阵是阶梯形矩阵的例子 (用双下划线标记阶梯点):
\[
\left(\begin{array}{cccccc}
\underline{1} & 2 & -1 & 0 & 7 & 10\\
0 & 0 & \underline{1} & -2 & 2 & 0\\
0 & 0 & 0 & \underline{5} & -1 & 1\\
0 & 0 & 0 & 0 & 0 & \underline{-1}\\
0 & 0 & 0 & 0 & 0 & 0
\end{array}\right),\quad
\left(\begin{array}{ccccc}
\underline{-11} & 20 & -\sqrt{2} & 5 & 7\\
0 & \underline{1} & 1 & -2 & 2\\
0 & 0 & 0 & \underline{5} & -1\\
0 & 0 & 0 & 0 & 0
\end{array}\right),
\]
\[
\left(\begin{array}{cccccc}
\underline{31} & 21 & -1 & 0 & 7 & 10\\
0 & 0 & \underline{1} & -2 & 2 & 0\\
0 & 0 & 0 & \underline{15} & 1 & 0
\end{array}\right),\quad
\left(\begin{array}{ccc}
\underline{3} & 2 & -1\\
0 & 0 & \underline{1}\\
0 & 0 & 0\\
0 & 0 & 0\\
0 & 0 & 0
\end{array}\right).
\]

下面几个矩阵不是阶梯形矩阵 (用双下划线标记阶梯点):
\[
\left(\begin{array}{cccccc}
\underline{1} & 0 & 0 & 0 & 0 & 0\\
\underline{2} & 1 & 0 & 0 & 0 & 0\\
\underline{1} & 2 & 3 & 1 & 0 & 0\\
\underline{-3} & -1 & 1 & 1 & 2 & 1
\end{array}\right),\quad
\left(\begin{array}{cccc}
\underline{1} & 2 & -1 & 0\\
0 & 0 & \underline{1} & -2\\
0 & 0 & \underline{2} & 5\\
0 & 0 & 0 & \underline{-1}
\end{array}\right),
\]
\[
\left(\begin{array}{cccccc}
\underline{1} & 2 & -1 & 0 & 7 & 10\\
\underline{-1} & 0 & 1 & -2 & 2 & 0\\
0 & 0 & 0 & \underline{5} & -1 & 1\\
0 & 0 & 0 & 0 & 0 & 0
\end{array}\right),\quad
\left(\begin{array}{cccccc}
\underline{1} & 2 & -1 & 0 & 7 & 10\\
0 & 0 & \underline{1} & -2 & 2 & 0\\
0 & 0 & 0 & \underline{5} & -1 & 1\\
0 & 0 & 0 & 0 & 0 & \underline{-1}\\
0 & 0 & 0 & 0 & 0 & \underline{2}
\end{array}\right).
\]

\noindent\textbf{定理 2.4.2}\quad 设 $A$ 是一个 $m\times n$ 矩阵，则经过若干次初等行变换，$A$ 可以化为阶梯形矩阵.

\noindent\textbf{证明}\quad 若 $A$ 的第一列元素全为零，则初等变换从第二列开始. 现设 $A$ 的第一列元素不全为零，则用行对换将非零元素换到第 $(1,1)$ 位置，再用第三类初等行变换即可将第一列其余元素消为零. 再看第二列，如果这时从第 $(2,2)$ 位置 (包括第 $(2,2)$ 元素) 以下全为零，则移至下一列. 若否，用行对换将非零元素换到第 $(2,2)$ 位置，再用第三类初等行变换消去这一列第 $(2,2)$ 位置以下的元素. 这样不断做下去，最后可以得到一个矩阵，其阶梯点的列指标随着行数严格递增，从而是阶梯形矩阵. $\square$

\noindent\textbf{例 2.4.3}\quad 用初等行变换化下列矩阵为阶梯形矩阵:
\[
\left(\begin{array}{ccc}
0 & 2 & -4\\
-1 & -4 & 5\\
3 & 1 & 7\\
0 & 5 & -10\\
2 & 3 & 0
\end{array}\right).
\]

\noindent\textbf{解}\quad
\[
\left(\begin{array}{ccc}
0 & 2 & -4\\
-1 & -4 & 5\\
3 & 1 & 7\\
0 & 5 & -10\\
2 & 3 & 0
\end{array}\right)
\longrightarrow
\left(\begin{array}{ccc}
-1 & -4 & 5\\
0 & 2 & -4\\
3 & 1 & 7\\
0 & 5 & -10\\
2 & 3 & 0
\end{array}\right)
\longrightarrow
\left(\begin{array}{ccc}
-1 & -4 & 5\\
0 & 2 & -4\\
0 & -11 & 22\\
0 & 5 & -10\\
0 & -5 & 10
\end{array}\right)
\longrightarrow
\left(\begin{array}{ccc}
-1 & -4 & 5\\
0 & 2 & -4\\
0 & 0 & 0\\
0 & 0 & 0\\
0 & 0 & 0
\end{array}\right).
\]

\subsection*{二、初等矩阵}

矩阵的初等变换能否通过矩阵的运算来实现？为此我们引进初等矩阵的概念.

\noindent\textbf{定义 2.4.4}\quad 对单位阵 $I_n$ 施以第一类、第二类、第三类初等变换后得到的矩阵分别称为第一类、第二类及第三类初等矩阵.

3 类初等矩阵的形状如下.

\noindent\textbf{第一类初等矩阵}\quad 第一类初等矩阵 $P_{ij}$ 表示将单位阵的第 $i$ 行与第 $j$ 行 (第 $i$ 列与第 $j$ 列) 对换后得到的矩阵:
\[
P_{ij}=\left(\begin{array}{ccccc}
1 & & & &\\
& \ddots & & &\\
& & 0 & \cdots & 1\\
& & \vdots & & \vdots\\
& & 1 & \cdots & 0\\
& & & \ddots &\\
& & & & 1
\end{array}\right).
\]

\noindent\textbf{第二类初等矩阵}\quad 第二类初等矩阵 $P_i(c)$ 表示将常数 $c(c\ne 0)$ 乘以单位阵的第 $i$ 行 (第 $i$ 列) 而得到的矩阵:
\[
P_i(c)=\left(\begin{array}{ccccc}
1 & & & &\\
& \ddots & & &\\
& & c & &\\
& & & \ddots &\\
& & & & 1
\end{array}\right).
\]

\noindent\textbf{第三类初等矩阵}\quad 第三类初等矩阵 $T_{ij}(c)$ 表示将单位阵的第 $i$ 行 (第 $j$ 列) 乘以 $c$ 后加到第 $j$ 行 (第 $i$ 列) 上得到的矩阵:
\[
T_{ij}(c)=\left(\begin{array}{cccccc}
1 & & & & &\\
& \ddots & & & &\\
& & 1 & \cdots & 0 &\\
& & \vdots & & \vdots &\\
& & c & \cdots & 1 &\\
& & & & & \ddots\\
& & & & & 1
\end{array}\right).
\]

下面的定理揭示了初等变换与初等矩阵的密切关系.

\noindent\textbf{定理 2.4.3}\quad 设 $A$ 是一个 $m\times n$ 阵，则对 $A$ 作一次初等行变换后得到的矩阵等于用一个 $m$ 阶相应的初等矩阵 (即第一类初等变换相应于第一类初等矩阵，第二类初等变换相应于第二类初等矩阵，等等) 左乘 $A$ 后得到的积. 矩阵 $A$ 作一次初等列变换后得到的矩阵等于用一个 $n$ 阶相应的初等矩阵右乘 $A$ 后所得的积.

\noindent\textbf{证明}\quad 我们只对初等行变换进行证明，对列的证明类似可得. 设 $A=(a_{ij})_{m\times n}$，由实际计算得
\[
P_{ij}A=
\left(\begin{array}{cccc}
a_{11} & a_{12} & \cdots & a_{1n}\\
\vdots & \vdots & & \vdots\\
a_{j1} & a_{j2} & \cdots & a_{jn}\\
\vdots & \vdots & & \vdots\\
a_{i1} & a_{i2} & \cdots & a_{in}\\
\vdots & \vdots & & \vdots\\
a_{m1} & a_{m2} & \cdots & a_{mn}
\end{array}\right),
\]
这等于将 $A$ 的第 $i$ 行与第 $j$ 行对换.
\[
P_i(c)A=
\left(\begin{array}{cccc}
a_{11} & a_{12} & \cdots & a_{1n}\\
\vdots & \vdots & & \vdots\\
ca_{i1} & ca_{i2} & \cdots & ca_{in}\\
\vdots & \vdots & & \vdots\\
a_{m1} & a_{m2} & \cdots & a_{mn}
\end{array}\right),
\]
这等于将 $A$ 的第 $i$ 行乘以 $c$. 而
\[
T_{ij}(c)A=
\left(\begin{array}{cccc}
a_{11} & a_{12} & \cdots & a_{1n}\\
\vdots & \vdots & & \vdots\\
a_{i1} & a_{i2} & \cdots & a_{in}\\
\vdots & \vdots & & \vdots\\
ca_{i1}+a_{j1} & ca_{i2}+a_{j2} & \cdots & ca_{in}+a_{jn}\\
\vdots & \vdots & & \vdots\\
a_{m1} & a_{m2} & \cdots & a_{mn}
\end{array}\right),
\]
这等于将 $A$ 的第 $i$ 行乘以 $c$ 后加到第 $j$ 行上去. $\square$

\noindent\textbf{注}\quad 我们通常用“行左列右”来表示定理中初等矩阵与初等变换的关系.

\noindent\textbf{推论 2.4.1}\quad 初等矩阵都是非异阵且其逆阵仍是同类初等矩阵:
\[
P_{ij}^{-1}=P_{ij},\quad P_i(c)^{-1}=P_i\left(\frac{1}{c}\right),\quad T_{ij}(c)^{-1}=T_{ij}(-c).
\]

\noindent\textbf{证明}\quad 由定理立即可得. $\square$

\noindent\textbf{推论 2.4.2}\quad 非异阵经初等变换后仍是非异阵，奇异阵经初等变换后仍是奇异阵.

\noindent\textbf{证明}\quad 因为初等矩阵都是可逆阵，可逆阵和可逆阵之积仍可逆，所以第一个结论成立. 又设 $A$ 是奇异阵，$P$ 是初等矩阵，假设 $PA$ 是非异阵，注意到 $P^{-1}$ 也是初等矩阵，故 $A=P^{-1}(PA)$ 将是非异阵，矛盾，因此 $PA$ 必是奇异阵. 同理 $AP$ 也是奇异阵. $\square$

\noindent\textbf{推论 2.4.3}\quad 3 类初等矩阵的行列式如下:
\[
|P_{ij}|=-1,\quad |P_i(c)|=c,\quad |T_{ij}(c)|=1.
\]

\noindent\textbf{证明}\quad 因为单位阵的行列式等于 1，故交换单位阵两行而得到的矩阵 $P_{ij}$ 的行列式等于 $-1$. 其余结论显然. $\square$

\noindent\textbf{定理 2.4.4}\quad 矩阵的相抵关系适合下列性质:

(1) $A\sim A$;

(2) 若 $A\sim B$，则 $B\sim A$;

(3) 若 $A\sim B,\ B\sim C$，则 $A\sim C$.

\noindent\textbf{证明}\quad (1) 显然单位阵 $I$ 也是初等矩阵，而 $IA=A$ 表明 $A\sim A$.

(2) 由 $A\sim B$ 知存在初等矩阵 $P_1,\cdots,P_m$ 以及初等矩阵 $Q_1,\cdots,Q_k$，使
\[
P_m\cdots P_1AQ_1\cdots Q_k=B.
\]
于是
\[
P_1^{-1}\cdots P_m^{-1}BQ_k^{-1}\cdots Q_1^{-1}=A.
\]
但是初等矩阵的逆阵仍是初等矩阵，因此 $B\sim A$.

(3) 由 $A\sim B,\ B\sim C$，可设
\[
B=P_m\cdots P_1AQ_1\cdots Q_k,\quad C=S_s\cdots S_1BT_1\cdots T_t,
\]
其中 $P_i,Q_i,S_i,T_i$ 都是初等矩阵. 于是
\[
C=S_s\cdots S_1P_m\cdots P_1AQ_1\cdots Q_kT_1\cdots T_t.
\]
因此 $A\sim C$. $\square$

\subsection*{习题 2.4}

1. 用初等变换将下列矩阵化为相抵标准型:
\[
(1)\quad
\left(\begin{array}{rrrr}
-1 & 0 & 1 & 2\\
3 & 1 & 0 & -1\\
0 & 2 & 1 & 4
\end{array}\right);
\qquad
(2)\quad
\left(\begin{array}{rrrr}
1 & 2 & 3 & 4\\
0 & -1 & 0 & -2\\
1 & 1 & 3 & 2\\
2 & 2 & 6 & 4
\end{array}\right).
\]

2. 用初等行变换将下列矩阵化为阶梯形矩阵:
\[
(1)\quad
\left(\begin{array}{rrr}
1 & 2 & 3\\
-1 & 0 & 1\\
0 & 1 & 2\\
2 & 2 & 4
\end{array}\right);
\qquad
(2)\quad
\left(\begin{array}{rrrrr}
1 & 0 & -1 & 2 & 3\\
2 & 2 & 1 & 3 & 8\\
3 & 4 & 4 & 5 & 11\\
1 & 2 & 4 & 3 & 1
\end{array}\right).
\]

3. 判断下列矩阵是否有相同的相抵标准型:
\[
(1)\quad
\left(\begin{array}{rrr}
1 & 2 & 3\\
2 & 4 & 6\\
0 & 1 & 2
\end{array}\right),\quad
\left(\begin{array}{rrr}
1 & -1 & 5\\
0 & 3 & 3\\
0 & 2 & 2
\end{array}\right);
\]
\[
(2)\quad
\left(\begin{array}{rrr}
-1 & 0 & 4\\
3 & 0 & -1\\
0 & 1 & -1
\end{array}\right),\quad
\left(\begin{array}{rrr}
1 & -1 & 5\\
-1 & 4 & -2\\
0 & 3 & 3
\end{array}\right).
\]

4. 求证：$n$ 阶方阵 $A$ 非异的充分必要条件是它和 $I_n$ 相抵.

5. 求证：对任意的 $m\times n$ 矩阵 $A$，总存在 $m$ 阶可逆阵 $P$ 和 $n$ 阶可逆阵 $Q$，使得 $PAQ$ 是相抵标准型.

\section*{§2.5\quad 矩阵乘积的行列式与初等变换法求逆阵}

\subsection*{一、矩阵乘积的行列式}

我们在上一节中证明了任意一个矩阵 $A$ 都可以通过初等变换化为相抵标准型，还证明了如果限制只用行变换，则可以将矩阵化为阶梯形. 如果矩阵 $A$ 是一个 $n$ 阶可逆阵，是否能够只用行变换就将它化为标准型呢？

\noindent\textbf{引理 2.5.1}\quad 设 $A$ 是一个 $n$ 阶可逆阵 (即非异阵)，则仅用初等行变换或仅用初等列变换就可以将它化为单位阵 $I_n$.

\noindent\textbf{证明}\quad 因为 $A$ 可逆，所以它没有整行或整列元素全为零，于是 $A$ 的第一列至少有一个非零元素，通过初等行变换可以将它换到第 $(1,1)$ 位置. 用这个非零元素经过第三类初等行变换就可以将第一列元素 (除第 $(1,1)$ 元素外) 全化为零. 于是不妨假设
\[
A=\left(\begin{array}{cccc}
a_{11} & a_{12} & \cdots & a_{1n}\\
0 & a_{22} & \cdots & a_{2n}\\
\vdots & \vdots & & \vdots\\
0 & a_{n2} & \cdots & a_{nn}
\end{array}\right).
\]
现在要说明 $a_{22},\cdots,a_{n2}$ 不全为零. 若否，则可以通过第三类初等列变换用 $a_{11}$ 消去 $a_{12}$，于是非异阵 $A$ 经过初等变换可化为一个第二列全为零的矩阵，即化为一个奇异阵，这与推论 2.4.2 矛盾. 既然 $a_{22},\cdots,a_{n2}$ 不全为零，我们又可以通过初等行变换将非零元素换到第 $(2,2)$ 位置. 再用这个非零元素经过第三类初等行变换就可以将第二列元素 (除第 $(2,2)$ 元素外) 全化为零. 不断这样做下去，即可将 $A$ 化为一个主对角元素全不为零的对角阵. 最后用第二类初等行变换即可将之化为单位阵 $I_n$. 同理可证仅用初等列变换也可以将非异阵 $A$ 化为单位阵 $I_n$. $\square$

\noindent\textbf{推论 2.5.1}\quad 任一 $n$ 阶非异阵均可表示成有限个初等矩阵的积.

\noindent\textbf{证明}\quad 由引理 2.5.1 知道，存在有限个初等矩阵 $P_1,\cdots,P_m$，使得
\[
P_m\cdots P_1A=I_n,
\]
因此
\[
A=P_1^{-1}\cdots P_m^{-1}.
\]
而初等矩阵的逆阵仍是初等矩阵，故结论成立. $\square$

下面我们要讨论这样一个问题：若 $A,B$ 都是 $n$ 阶矩阵，它们积的行列式 $|AB|$ 和 $|A|,|B|$ 有什么关系？我们先讨论简单的情形，即 $A,B$ 中至少有一个是初等矩阵的情形. 由行列式性质，很容易得到下面的引理.

\noindent\textbf{引理 2.5.2}\quad 设 $A$ 是一个 $n$ 阶方阵，$Q$ 是一个 $n$ 阶初等矩阵，则
\[
|QA|=|Q||A|=|AQ|.
\]

\noindent\textbf{证明}\quad 若 $Q=P_{ij}$，则 $QA$ 为 $A$ 的第 $i$ 行及第 $j$ 行对换后所得之矩阵，其行列式值等于 $-|A|$. 又 $|P_{ij}|=-1$，因此 $|P_{ij}A|=-|A|=|P_{ij}||A|$. 若 $Q=P_i(c)(c\ne 0)$，不难验证 $|P_i(c)A|=c|A|=|P_i(c)||A|$. 若 $Q=T_{ij}(c)$，同样不难验证 $|T_{ij}(c)A|=|A|=|T_{ij}(c)||A|$. 当 $Q$ 作用在 $A$ 的右侧时也可类似证明. $\square$

\noindent\textbf{定理 2.5.1}\quad 一个 $n$ 阶方阵 $A$ 为非异阵的充分必要条件是它的行列式的值不等于零.

\noindent\textbf{证明}\quad 由定理 2.3.1 知道若 $|A|\ne 0$，则 $A$ 是非异阵，故只需证明若 $A$ 非异，则 $|A|\ne 0$. 由推论 2.5.1，非异阵 $A$ 等于有限个初等矩阵之积，设 $A=P_1P_2\cdots P_t$，则从上面的引理知道，$|A|=|P_1||P_2|\cdots |P_t|$. 由于初等矩阵的行列式不等于零，故 $|A|\ne 0$. $\square$

现在我们可以证明下述重要的行列式乘法定理.

\noindent\textbf{定理 2.5.2}\quad 设 $A,B$ 都是 $n$ 阶矩阵，则
\[
|AB|=|A||B|.
\]

\noindent\textbf{证明}\quad 分两种情形进行讨论. 第一种情形，设 $A$ 是非异阵，则存在若干个初等矩阵 $Q_1,\cdots,Q_m$，使得 $A=Q_1\cdots Q_m$，故
\[
|AB|=|Q_1\cdots Q_mB|=|Q_1|\cdots |Q_m||B|=|Q_1\cdots Q_m||B|=|A||B|.
\]
第二种情形，设 $A$ 为奇异阵，这时 $|A|=0$，故只需证明 $|AB|=0$. 因为 $A$ 是奇异阵，故存在初等矩阵 $P_1,\cdots,P_s;Q_1,\cdots,Q_r$，使得
\[
P_s\cdots P_1AQ_1\cdots Q_r=D,
\]
其中 $D$ 是 $A$ 的相抵标准型，它是一个对角阵，主对角元素为 1 或 0. 因为 $A$ 奇异，所以 $D$ 也是奇异阵，于是至少最后一行元素全为零. 注意到
\[
P_s\cdots P_1A=DQ_r^{-1}\cdots Q_1^{-1},
\]
由矩阵乘法可知，因为 $D$ 的第 $n$ 行元素全为零，所以 $DQ_r^{-1}\cdots Q_1^{-1}=P_s\cdots P_1A$ 的第 $n$ 行元素全为零，从而 $P_s\cdots P_1AB$ 的第 $n$ 行元素也全为零. 于是
\[
|P_s|\cdots |P_1||AB|=|P_s\cdots P_1AB|=0,
\]
又 $|P_i|\ne 0(i=1,\cdots,s)$，故 $|AB|=0=|A||B|$. $\square$

\noindent\textbf{推论 2.5.2}\quad 一个奇异阵与任一同阶方阵之积仍为奇异阵，两个同阶非异阵之积仍为非异阵.

\noindent\textbf{证明}\quad 由定理 2.5.2 及定理 2.5.1 即得. $\square$

\noindent\textbf{推论 2.5.3}\quad 若 $A$ 是非异阵，则 $|A^{-1}|=|A|^{-1}$.

\noindent\textbf{证明}\quad 由 $1=|I_n|=|AA^{-1}|=|A||A^{-1}|$ 即得 $|A^{-1}|=|A|^{-1}$. $\square$

\noindent\textbf{推论 2.5.4}\quad 若 $A,B$ 都是 $n$ 阶方阵且 $AB=I_n$ (或 $BA=I_n$)，则 $BA=I_n$ (或 $AB=I_n$)，即 $B=A^{-1}$.

\noindent\textbf{证明}\quad 因为 $|A||B|=|AB|=|I_n|=1$，所以 $|A|\ne 0$，即 $A$ 是非异阵. 设 $C$ 是 $A$ 的逆阵，则 $CA=I_n$，而 $B=I_nB=(CA)B=C(AB)=CI_n=C$，因此 $B=A^{-1}$. 同理若 $BA=I_n$，也可推出 $AB=I_n$. $\square$

\noindent\textbf{例 2.5.1}\quad 计算下列 $n+1$ 阶矩阵 $A$ 的行列式:
\[
A=\left(\begin{array}{cccc}
(a_0+b_0)^n & (a_0+b_1)^n & \cdots & (a_0+b_n)^n\\
(a_1+b_0)^n & (a_1+b_1)^n & \cdots & (a_1+b_n)^n\\
\vdots & \vdots & & \vdots\\
(a_n+b_0)^n & (a_n+b_1)^n & \cdots & (a_n+b_n)^n
\end{array}\right).
\]

\noindent\textbf{解}\quad 将 $A$ 分解为两个矩阵之积:
\[
A=
\left(\begin{array}{ccccc}
1 & C_n^1a_0 & C_n^2a_0^2 & \cdots & C_n^na_0^n\\
1 & C_n^1a_1 & C_n^2a_1^2 & \cdots & C_n^na_1^n\\
\vdots & \vdots & \vdots & & \vdots\\
1 & C_n^1a_n & C_n^2a_n^2 & \cdots & C_n^na_n^n
\end{array}\right)
\left(\begin{array}{ccccc}
b_0^n & b_1^n & b_2^n & \cdots & b_n^n\\
b_0^{n-1} & b_1^{n-1} & b_2^{n-1} & \cdots & b_n^{n-1}\\
\vdots & \vdots & \vdots & & \vdots\\
1 & 1 & 1 & \cdots & 1
\end{array}\right).
\]
上式左边矩阵的行列式每一列提出公因子后就是一个 Vandermonde 行列式. 右边矩阵的行列式也可以化为 Vandermonde 行列式并求出其值，于是
\[
|A|=C_n^1C_n^2\cdots C_n^n\prod_{0\le i<j\le n}(a_j-a_i)(b_i-b_j).
\]

\noindent\textbf{例 2.5.2}\quad 计算行列式:
\[
\left|\begin{array}{rrrr}
x & -y & -z & -w\\
y & x & -w & z\\
z & w & x & -y\\
w & -z & y & x
\end{array}\right|.
\]

\noindent\textbf{解}\quad 设该行列式代表的矩阵为 $A$，则
\[
AA'=
\left(\begin{array}{rrrr}
x & -y & -z & -w\\
y & x & -w & z\\
z & w & x & -y\\
w & -z & y & x
\end{array}\right)
\left(\begin{array}{rrrr}
x & y & z & w\\
-y & x & w & -z\\
-z & -w & x & y\\
-w & z & -y & x
\end{array}\right)
=
\left(\begin{array}{cccc}
u & 0 & 0 & 0\\
0 & u & 0 & 0\\
0 & 0 & u & 0\\
0 & 0 & 0 & u
\end{array}\right),
\]
其中 $u=x^2+y^2+z^2+w^2$. 因此
\[
|A|^2=(x^2+y^2+z^2+w^2)^4.
\]
令 $x=1,y=z=w=0$，显然 $|A|=1$，故
\[
|A|=(x^2+y^2+z^2+w^2)^2.
\]

\subsection*{二、初等变换法求逆阵}

我们知道，用伴随阵求非异阵的逆阵是非常麻烦的，有没有更加简单的方法？由引理 2.5.1 知道，任意一个可逆阵都可以只用初等行变换将它化为单位阵. 若 $A$ 是可逆阵，则存在初等矩阵 $Q_1,Q_2,\cdots,Q_t$，使得
\[
Q_1Q_2\cdots Q_tA=I_n,
\]
故
\[
A^{-1}=Q_1Q_2\cdots Q_t=Q_1Q_2\cdots Q_tI_n.
\]
这两个式子启发我们可以这样来求逆阵:

作一个 $n\times 2n$ 矩阵 $(A\vdots I_n)$，这个矩阵的前 $n$ 列为 $A$，后 $n$ 列为单位阵 $I_n$. 对矩阵 $(A\vdots I_n)$ 进行初等行变换把 $A$ 变成 $I_n$，这时右边的 $I_n$ 就变成了 $A^{-1}$. 我们下面举例来说明这个方法.

\noindent\textbf{例 2.5.3}\quad 求下列非异阵 $A$ 的逆阵:
\[
A=\left(\begin{array}{ccc}
1 & 2 & 3\\
2 & 1 & 2\\
1 & 3 & 4
\end{array}\right).
\]

\noindent\textbf{解}\quad
\[
(A\vdots I_3)=
\left(\begin{array}{ccc|ccc}
1 & 2 & 3 & 1 & 0 & 0\\
2 & 1 & 2 & 0 & 1 & 0\\
1 & 3 & 4 & 0 & 0 & 1
\end{array}\right).
\]
对 $(A\vdots I_3)$ 进行初等行变换:
\[
\left(\begin{array}{ccc|ccc}
1 & 2 & 3 & 1 & 0 & 0\\
2 & 1 & 2 & 0 & 1 & 0\\
1 & 3 & 4 & 0 & 0 & 1
\end{array}\right)
\longrightarrow
\left(\begin{array}{ccc|ccc}
1 & 2 & 3 & 1 & 0 & 0\\
0 & -3 & -4 & -2 & 1 & 0\\
0 & 1 & 1 & -1 & 0 & 1
\end{array}\right)
\longrightarrow
\left(\begin{array}{ccc|ccc}
1 & 2 & 3 & 1 & 0 & 0\\
0 & 1 & 1 & -1 & 0 & 1\\
0 & -3 & -4 & -2 & 1 & 0
\end{array}\right)
\]
\[
\longrightarrow
\left(\begin{array}{ccc|ccc}
1 & 0 & 1 & 3 & 0 & -2\\
0 & 1 & 1 & -1 & 0 & 1\\
0 & 0 & -1 & -5 & 1 & 3
\end{array}\right)
\longrightarrow
\left(\begin{array}{ccc|ccc}
1 & 0 & 0 & -2 & 1 & 1\\
0 & 1 & 0 & -6 & 1 & 4\\
0 & 0 & -1 & -5 & 1 & 3
\end{array}\right)
\longrightarrow
\left(\begin{array}{ccc|ccc}
1 & 0 & 0 & -2 & 1 & 1\\
0 & 1 & 0 & -6 & 1 & 4\\
0 & 0 & 1 & 5 & -1 & -3
\end{array}\right).
\]
于是
\[
A^{-1}=\left(\begin{array}{rrr}
-2 & 1 & 1\\
-6 & 1 & 4\\
5 & -1 & -3
\end{array}\right).
\]

\noindent\textbf{注}\quad 在用初等变换法求逆阵的整个过程中，对 $(A\vdots I_n)$ 只能用初等行变换而不能用初等列变换. 请读者考虑这样一个问题：

如果我们只用初等列变换，如何来求已知可逆阵的逆阵？

\noindent\textbf{例 2.5.4}\quad 求解下列矩阵方程:
\[
\left(\begin{array}{rrr}
1 & 0 & 1\\
-1 & 1 & 1\\
2 & -1 & 1
\end{array}\right)X=
\left(\begin{array}{rr}
1 & 1\\
0 & 1\\
-1 & 0
\end{array}\right).
\]

\noindent\textbf{解}\quad 因为
\[
\left|\begin{array}{rrr}
1 & 0 & 1\\
-1 & 1 & 1\\
2 & -1 & 1
\end{array}\right|\ne 0,
\]
故
\[
X=
\left(\begin{array}{rrr}
1 & 0 & 1\\
-1 & 1 & 1\\
2 & -1 & 1
\end{array}\right)^{-1}
\left(\begin{array}{rr}
1 & 1\\
0 & 1\\
-1 & 0
\end{array}\right).
\]
我们用与求逆阵类似的初等变换法求 $X$:
\[
\left(\begin{array}{rrr|rr}
1 & 0 & 1 & 1 & 1\\
-1 & 1 & 1 & 0 & 1\\
2 & -1 & 1 & -1 & 0
\end{array}\right)
\longrightarrow
\left(\begin{array}{rrr|rr}
1 & 0 & 1 & 1 & 1\\
0 & 1 & 2 & 1 & 2\\
0 & -1 & -1 & -3 & -2
\end{array}\right)
\]
\[
\longrightarrow
\left(\begin{array}{rrr|rr}
1 & 0 & 1 & 1 & 1\\
0 & 1 & 2 & 1 & 2\\
0 & 0 & 1 & -2 & 0
\end{array}\right)
\longrightarrow
\left(\begin{array}{rrr|rr}
1 & 0 & 0 & 3 & 1\\
0 & 1 & 0 & 5 & 2\\
0 & 0 & 1 & -2 & 0
\end{array}\right).
\]
因此
\[
X=\left(\begin{array}{rr}
3 & 1\\
5 & 2\\
-2 & 0
\end{array}\right).
\]

\subsection*{习题 2.5}

1. 用初等变换法求下列矩阵的逆阵:
\[
(1)\quad
\left(\begin{array}{rrrr}
1 & 2 & 3 & 4\\
2 & 3 & 1 & 2\\
1 & 1 & -1 & -1\\
1 & 0 & -6 & -6
\end{array}\right);
\qquad
(2)\quad
\left(\begin{array}{rrr}
1 & 1 & 0\\
-2 & 1 & -1\\
3 & 1 & 1
\end{array}\right);
\qquad
(3)\quad
\left(\begin{array}{rrr}
1 & 2 & -1\\
3 & 1 & 0\\
7 & 6 & -2
\end{array}\right).
\]

2. 求下列 $n$ 阶方阵的逆阵:
\[
A=\left(\begin{array}{ccccc}
0 & a_1 & 0 & \cdots & 0\\
0 & 0 & a_2 & \cdots & 0\\
\vdots & \vdots & \vdots & & \vdots\\
0 & 0 & 0 & \cdots & a_{n-1}\\
a_n & 0 & 0 & \cdots & 0
\end{array}\right),
\]
其中 $a_i\ne 0(i=1,2,\cdots,n)$.

3. 求下列矩阵方程的解:
\[
\left(\begin{array}{rrr}
2 & 2 & 3\\
1 & -1 & 0\\
-1 & 2 & 1
\end{array}\right)X=
\left(\begin{array}{rr}
4 & -1\\
2 & 1\\
1 & 0
\end{array}\right).
\]

4. 求下列矩阵方程的解:
\[
X\left(\begin{array}{rrr}
2 & 2 & 3\\
1 & -1 & 0\\
-1 & 2 & 1
\end{array}\right)=
\left(\begin{array}{rrr}
1 & 1 & 2\\
0 & -1 & 3
\end{array}\right).
\]

5. 设
\[
S=\left(\begin{array}{ccccc}
s_0 & s_1 & s_2 & \cdots & s_{n-1}\\
s_1 & s_2 & s_3 & \cdots & s_n\\
s_2 & s_3 & s_4 & \cdots & s_{n+1}\\
\vdots & \vdots & \vdots & & \vdots\\
s_{n-1} & s_n & s_{n+1} & \cdots & s_{2n-2}
\end{array}\right),
\]
其中 $s_k=x_1^k+x_2^k+\cdots+x_n^k$. 计算 $|S|$ 并证明 $|S|\ge 0$ 对一切实数 $x_i(i=1,2,\cdots,n)$ 成立.

6. 利用行列式乘法证明循环矩阵的行列式的值等于 $f(\varepsilon_1)f(\varepsilon_2)\cdots f(\varepsilon_n)$，其中
\[
f(x)=a_1+a_2x+a_3x^2+\cdots+a_nx^{n-1},
\]
$\varepsilon_i(i=1,2,\cdots,n)$ 为 1 的全体 $n$ 次方根:
\[
\left|\begin{array}{ccccc}
a_1 & a_2 & a_3 & \cdots & a_n\\
a_n & a_1 & a_2 & \cdots & a_{n-1}\\
a_{n-1} & a_n & a_1 & \cdots & a_{n-2}\\
\vdots & \vdots & \vdots & & \vdots\\
a_2 & a_3 & a_4 & \cdots & a_1
\end{array}\right|.
\]

7. 利用习题 6 的结论计算复习题一第 5 题.

\section*{§2.6\quad 分块矩阵}

矩阵运算是一种比较复杂的运算. 为了简化这种运算，我们引进分块矩阵及其运算的概念. 请读者务必注意，分块矩阵及其运算不是新的运算，而是矩阵运算的简化形式.

什么叫矩阵的分块？简单地说就是用横虚线与竖虚线将一个矩阵分成若干块，这样得到的矩阵就称为“分块矩阵”. 例如:
\[
A=\left(\begin{array}{ccc|c}
1 & 2 & -1 & 0\\
2 & 5 & 0 & -2\\ \hline
3 & 1 & -1 & 3
\end{array}\right)
\]
是一个分块矩阵，若记
\[
A_{11}=\left(\begin{array}{rrr}
1 & 2 & -1\\
2 & 5 & 0
\end{array}\right),\quad
A_{12}=\left(\begin{array}{r}
0\\
-2
\end{array}\right);
\]
\[
A_{21}=(3,1,-1),\quad A_{22}=(3),
\]
则 $A$ 可表示为
\[
A=\left(\begin{array}{cc}
A_{11} & A_{12}\\
A_{21} & A_{22}
\end{array}\right),
\]
这是一个分了 4 块的矩阵.

一般地，对 $m\times n$ 矩阵 $A$，若先用若干条横虚线把它分成 $r$ 块，再用若干条竖虚线把它分成 $s$ 块，我们就得到了一个 $rs$ 块的分块矩阵，可记为
\[
A=\left(\begin{array}{cccc}
A_{11} & A_{12} & \cdots & A_{1s}\\
A_{21} & A_{22} & \cdots & A_{2s}\\
\vdots & \vdots & & \vdots\\
A_{r1} & A_{r2} & \cdots & A_{rs}
\end{array}\right).
\]
注意，这里 $A_{ij}$ 代表一个矩阵，$A_{ij}$ 常称为 $A$ 的第 $(i,j)$ 块. $A$ 也可记为 $A=(A_{ij})$，但需注明这是分块矩阵.

一个矩阵可以有各种各样的分块方法，究竟怎么分比较好，要看具体需要而定. 例如:
\[
A=\left(\begin{array}{cc|cc|c}
1 & 1 & 0 & 0 & 0\\
-1 & 1 & 0 & 0 & 0\\ \hline
0 & 0 & 1 & 0 & 0\\
0 & 0 & 1 & 1 & 0\\ \hline
0 & 0 & 0 & 0 & 1
\end{array}\right).
\]
记
\[
A_1=\left(\begin{array}{rr}
1 & 1\\
-1 & 1
\end{array}\right),\quad
A_2=\left(\begin{array}{rr}
1 & 0\\
1 & 1
\end{array}\right),\quad A_3=(1),
\]
则
\[
A=\left(\begin{array}{ccc}
A_1 & O & O\\
O & A_2 & O\\
O & O & A_3
\end{array}\right).
\]
$A$ 作为分块矩阵看，是一个对角阵. 这种矩阵称为分块对角阵，尽管它本身并不是对角阵，需要注意的是 $A$ 中的 $O$ 都表示零矩阵.

两个分块矩阵 $A=(A_{ij})_{r\times s}$ 及 $B=(B_{ij})_{k\times l}$ 称为相等，若 $r=k,\ s=l$，且 $A_{ij}=B_{ij}(i=1,2,\cdots,r;\ j=1,2,\cdots,s)$. 因此两个分块矩阵相等，不仅它们的分块方式相同，而且每一块也相等. 显然，这时 $A$ 与 $B$ 作为普通矩阵也相等.

下面我们依次来研究分块矩阵的运算.

\subsection*{一、分块矩阵的加减法}

设有 $m\times n$ 矩阵 $A$ 及 $B$，它们具有相同的分块，即
\[
A=(A_{ij})_{r\times s},\quad B=(B_{ij})_{r\times s},
\]
且 $A_{ij}$ 与 $B_{ij}$ 作为矩阵其行数与列数分别相等，则
\[
A+B=(A_{ij}+B_{ij}),\quad A-B=(A_{ij}-B_{ij}).
\]
显然，两个分块矩阵之和 (或差) 仍是一个分块矩阵，且这个和 (或差) 与 $A,B$ 作为普通矩阵之和 (或差) 是一致的.

\subsection*{二、分块矩阵的数乘}

分块矩阵 $A=(A_{ij})_{r\times s}$ 与常数 $c$ 的数乘为
\[
cA=(cA_{ij})_{r\times s}.
\]

\subsection*{三、分块矩阵的乘法}

分块矩阵的乘法与普通矩阵的乘法在形式上类似，只是在处理块与块之间的乘法时必须保证符合矩阵相乘的条件. 因此，对分块的情况要特别予以注意. 设
\[
A=(A_{ij})_{r\times s},\quad B=(B_{ij})_{s\times t}
\]
是两个分块矩阵 (注意：$A$ 的列分成 $s$ 块而 $B$ 的行也分成 $s$ 块). 又设
\[
A=\left(\begin{array}{cccc}
A_{11} & A_{12} & \cdots & A_{1s}\\
A_{21} & A_{22} & \cdots & A_{2s}\\
\vdots & \vdots & & \vdots\\
A_{r1} & A_{r2} & \cdots & A_{rs}
\end{array}\right),\quad
B=\left(\begin{array}{cccc}
B_{11} & B_{12} & \cdots & B_{1t}\\
B_{21} & B_{22} & \cdots & B_{2t}\\
\vdots & \vdots & & \vdots\\
B_{s1} & B_{s2} & \cdots & B_{st}
\end{array}\right).
\]
上述分块矩阵适合如下条件：在 $A$ 中，第 $(1,1)$ 块 $A_{11}$ 的行数为 $m_1$，列数为 $n_1$，第 $(1,2)$ 块 $A_{12}$ 的行数为 $m_1$，列数为 $n_2,\cdots$，第 $(i,j)$ 块 $A_{ij}$ 的行数为 $m_i$，列数为 $n_j$. $B$ 中第 $(i,j)$ 块 $B_{ij}$ 的行数为 $n_i$，列数为 $l_j$. 这样的分块方式保证了分块相乘有意义. 若记分块矩阵 $A$ 与 $B$ 的积为
\[
C=\left(\begin{array}{cccc}
C_{11} & C_{12} & \cdots & C_{1t}\\
C_{21} & C_{22} & \cdots & C_{2t}\\
\vdots & \vdots & & \vdots\\
C_{r1} & C_{r2} & \cdots & C_{rt}
\end{array}\right),
\]
则 $C_{ij}$ 是一个 $m_i\times l_j$ 矩阵，且
\[
C_{ij}=A_{i1}B_{1j}+A_{i2}B_{2j}+\cdots+A_{is}B_{sj}.
\]

\noindent\textbf{例 2.6.1}\quad
\[
A=\left(\begin{array}{cc|c|cc}
1 & 0 & 2 & -1 & 0\\
0 & 1 & 1 & -2 & 1\\ \hline
0 & 0 & 3 & 1 & 0\\
1 & 0 & -2 & 0 & 1
\end{array}\right),\quad
B=\left(\begin{array}{cc|c}
1 & 0 & 2\\
0 & 1 & 0\\ \hline
-1 & 1 & 3\\ \hline
0 & 1 & -1\\
2 & 0 & 1
\end{array}\right),
\]
求 $AB$.

\noindent\textbf{解}\quad 将 $A,B$ 写成下列分块形状:
\[
A=\left(\begin{array}{ccc}
A_{11} & A_{12} & A_{13}\\
A_{21} & A_{22} & A_{23}
\end{array}\right),\quad
B=\left(\begin{array}{cc}
B_{11} & B_{12}\\
B_{21} & B_{22}\\
B_{31} & B_{32}
\end{array}\right),
\]
其中 $A_{ij},B_{ij}$ 是 $A,B$ 中相应的块. 容易看出这两个分块矩阵符合相乘的条件. 设 $C=AB$，则 $C$ 也是分块矩阵. 因为 $A$ 是 $2\times 3$ 分块，$B$ 是 $3\times 2$ 分块，故 $C$ 是 $2\times 2$ 分块. 不难看出
\[
C_{11}=A_{11}B_{11}+A_{12}B_{21}+A_{13}B_{31}.
\]
将各块代入，得到
\[
\begin{aligned}
C_{11}
&=\left(\begin{array}{cc}
1 & 0\\
0 & 1
\end{array}\right)
\left(\begin{array}{cc}
1 & 0\\
0 & 1
\end{array}\right)
+\left(\begin{array}{c}
2\\
1
\end{array}\right)(-1,1)
+\left(\begin{array}{rr}
-1 & 0\\
-2 & 1
\end{array}\right)
\left(\begin{array}{cc}
0 & 1\\
2 & 0
\end{array}\right)\\
&=\left(\begin{array}{cc}
1 & 0\\
0 & 1
\end{array}\right)
+\left(\begin{array}{rr}
-2 & 2\\
-1 & 1
\end{array}\right)
+\left(\begin{array}{rr}
0 & -1\\
2 & -2
\end{array}\right)
=\left(\begin{array}{rr}
-1 & 1\\
1 & 0
\end{array}\right).
\end{aligned}
\]
同理
\[
C_{12}=A_{11}B_{12}+A_{12}B_{22}+A_{13}B_{32}=\left(\begin{array}{c}
9\\
6
\end{array}\right),
\]
\[
C_{21}=A_{21}B_{11}+A_{22}B_{21}+A_{23}B_{31}=
\left(\begin{array}{rr}
-3 & 4\\
5 & -2
\end{array}\right),
\]
\[
C_{22}=A_{21}B_{12}+A_{22}B_{22}+A_{23}B_{32}=
\left(\begin{array}{r}
8\\
-3
\end{array}\right).
\]
于是
\[
C=\left(\begin{array}{rr|r}
-1 & 1 & 9\\
1 & 0 & 6\\ \hline
-3 & 4 & 8\\
5 & -2 & -3
\end{array}\right).
\]

读者也许感到，上述例子似乎比不分块更麻烦. 现在来看几个例子，它们表明了分块运算的优越性.

\noindent\textbf{例 2.6.2}\quad 设有两个分块对角阵:
\[
A=\left(\begin{array}{cccc}
A_1 & O & \cdots & O\\
O & A_2 & \cdots & O\\
\vdots & \vdots & & \vdots\\
O & O & \cdots & A_k
\end{array}\right),\quad
B=\left(\begin{array}{cccc}
B_1 & O & \cdots & O\\
O & B_2 & \cdots & O\\
\vdots & \vdots & & \vdots\\
O & O & \cdots & B_k
\end{array}\right),
\]
其中 $A_i$ 与 $B_i$ 都是同阶方阵，因此 $A$ 与 $B$ 可按分块矩阵相乘，
\[
AB=\left(\begin{array}{cccc}
A_1B_1 & O & \cdots & O\\
O & A_2B_2 & \cdots & O\\
\vdots & \vdots & & \vdots\\
O & O & \cdots & A_kB_k
\end{array}\right).
\]
这个例子表明，分块对角阵相乘时只需将主对角线上的块相乘即可.

\noindent\textbf{例 2.6.3}\quad $A$ 是一个分块对角阵:
\[
A=\left(\begin{array}{cccc}
A_1 & O & \cdots & O\\
O & A_2 & \cdots & O\\
\vdots & \vdots & & \vdots\\
O & O & \cdots & A_k
\end{array}\right),
\]
其中每块 $A_i$ 都是非异阵，求证：$A$ 也是非异阵.

\noindent\textbf{证明}\quad 设 $A_i$ 之逆为 $A_i^{-1}$，则显然
\[
A^{-1}=\left(\begin{array}{cccc}
A_1^{-1} & O & \cdots & O\\
O & A_2^{-1} & \cdots & O\\
\vdots & \vdots & & \vdots\\
O & O & \cdots & A_k^{-1}
\end{array}\right). \quad \square
\]

又如，设 $A$ 是一个 $m\times n$ 矩阵，$B$ 是一个 $n\times r$ 矩阵，可对 $B$ 作列分块，即将 $B$ 的每个列向量分作一块，记为 $\beta_j(j=1,2,\cdots,r)$，则
\[
B=(\beta_1,\beta_2,\cdots,\beta_r).
\]
又将 $A$ 看成是只分成一块的矩阵，则 $AB$ 可按分块矩阵相乘，且 $AB$ 的列分块为
\[
AB=(A\beta_1,A\beta_2,\cdots,A\beta_r).
\]
同样，可对 $A$ 作行分块，即将 $A$ 的每个行向量分作一块，记为 $\alpha_i(i=1,2,\cdots,m)$，则
\[
A=\left(\begin{array}{c}
\alpha_1\\
\alpha_2\\
\vdots\\
\alpha_m
\end{array}\right).
\]
又将 $B$ 看成是只有一块的矩阵，则 $AB$ 可按分块矩阵相乘，且 $AB$ 的行分块为
\[
AB=\left(\begin{array}{c}
\alpha_1B\\
\alpha_2B\\
\vdots\\
\alpha_mB
\end{array}\right).
\]

\subsection*{四、分块矩阵的转置}

设有分块矩阵 $A=(A_{ij})_{r\times s}$，则 $A$ 的转置为 $s\times r$ 分块矩阵:
\[
A'=\left(\begin{array}{cccc}
A'_{11} & A'_{21} & \cdots & A'_{r1}\\
A'_{12} & A'_{22} & \cdots & A'_{r2}\\
\vdots & \vdots & & \vdots\\
A'_{1s} & A'_{2s} & \cdots & A'_{rs}
\end{array}\right).
\]

\subsection*{五、分块矩阵的共轭}

设 $A$ 是一个分块复矩阵，$A=(A_{ij})_{r\times s}$，则 $A$ 的共轭矩阵也是一个 $r\times s$ 分块矩阵，且
\[
\overline{A}=\left(\begin{array}{cccc}
\overline{A}_{11} & \overline{A}_{12} & \cdots & \overline{A}_{1s}\\
\overline{A}_{21} & \overline{A}_{22} & \cdots & \overline{A}_{2s}\\
\vdots & \vdots & & \vdots\\
\overline{A}_{r1} & \overline{A}_{r2} & \cdots & \overline{A}_{rs}
\end{array}\right).
\]

对分块矩阵而言，也有分块初等变换和分块初等矩阵的概念. 它们是处理分块矩阵问题的有用工具.

分块初等变换和普通的初等变换类似，包含 3 类:

第一类：对调分块矩阵的两个分块行或两个分块列;

第二类：以某个可逆阵左乘以分块矩阵的某个分块行，或右乘以某个分块列;

第三类：以某个矩阵左乘以分块矩阵的某个分块行后加到另一分块行上去，或以某个矩阵右乘以分块矩阵的某个分块列后加到另一分块列上去.

我们假设上面所提到的运算都是可以进行的.

分块初等矩阵也和普通的初等矩阵类似.

记 $I=\operatorname{diag}\{I_{m_1},I_{m_2},\cdots,I_{m_k}\}$ 是分块单位阵，定义下列 3 种矩阵为 3 类分块初等矩阵:

第一类：对调 $I$ 的第 $i$ 分块行 (列) 与第 $j$ 分块行 (列) 得到的矩阵;

第二类：以可逆阵 $C$ 左 (右) 乘以 $I$ 的第 $i$ 分块行 (列) 得到的矩阵;

第三类：以矩阵 $B$ 左 (右) 乘以 $I$ 的第 $i$ 分块行 (列) 后加到第 $j$ 分块行 (列) 上得到的矩阵.

容易验证分块初等矩阵都是可逆阵，其中第三类分块初等矩阵的行列式的值等于 1，并且矩阵的分块初等行 (列) 变换相当于用同类分块初等矩阵左 (右) 乘以被变换的矩阵.

\noindent\textbf{注}\quad 由上面的结论知道，进行第三类分块初等变换不改变矩阵的行列式的值. 更一般地，进行分块初等变换不改变矩阵的秩 (秩的概念将在下一章介绍).

\noindent\textbf{引理 2.6.1}\quad 设 $A,C$ 分别是 $m,n$ 阶方阵，则对分块上 (下) 三角行列式有:
\[
|G|=\left|\begin{array}{cc}
A & B\\
O & C
\end{array}\right|=|A||C|,\quad
|H|=\left|\begin{array}{cc}
A & O\\
B & C
\end{array}\right|=|A||C|.
\]

\noindent\textbf{证明}\quad 本命题可以用 Laplace 定理立即得到，但我们采用另外一种方法来证明. 只对第一式进行证明，第二式同理可得. 对 $A$ 的阶 $m$ 用归纳法，当 $m=1$ 时，结论显然成立. 假设对左上角是 $m-1$ 阶矩阵的分块上三角行列式结论为真. 设
\[
A=\left(\begin{array}{cccc}
a_{11} & a_{12} & \cdots & a_{1m}\\
a_{21} & a_{22} & \cdots & a_{2m}\\
\vdots & \vdots & & \vdots\\
a_{m1} & a_{m2} & \cdots & a_{mm}
\end{array}\right),
\]
将 $|G|$ 按第一列展开 (注意到第一列从第 $m+1$ 行起的元素全为零):
\[
|G|=a_{11}G_{11}+a_{21}G_{21}+\cdots+a_{m1}G_{m1},
\]
其中 $G_{i1}$ 是 $a_{i1}$ 在 $|G|$ 中的代数余子式. 与每个 $G_{i1}$ 相应的余子式是一个左上角为 $m-1$ 阶矩阵的分块上三角行列式，由归纳假设可得
\[
G_{i1}=A_{i1}|C|,
\]
其中 $A_{i1}$ 是元素 $a_{i1}$ 在 $|A|$ 中的代数余子式. 因此
\[
|G|=a_{11}A_{11}|C|+a_{21}A_{21}|C|+\cdots+a_{m1}A_{m1}|C|=|A||C|. \quad \square
\]

\noindent\textbf{定理 2.6.1}\quad 若 $A$ 是 $m$ 阶可逆阵，$D$ 是 $n$ 阶矩阵，$B$ 为 $m\times n$ 矩阵，$C$ 为 $n\times m$ 矩阵，则
\[
\left|\begin{array}{cc}
A & B\\
C & D
\end{array}\right|=|A||D-CA^{-1}B|.
\]
若 $D$ 可逆 (这时 $A$ 不必假设可逆)，则有
\[
\left|\begin{array}{cc}
A & B\\
C & D
\end{array}\right|=|D||A-BD^{-1}C|.
\]

\noindent\textbf{证明}\quad 用第三类分块初等变换，以 $-CA^{-1}$ 左乘以第一分块行加到第二分块行上得到
\[
\left(\begin{array}{cc}
A & B\\
C & D
\end{array}\right)\longrightarrow
\left(\begin{array}{cc}
A & B\\
O & D-CA^{-1}B
\end{array}\right).
\]
第三类分块初等变换不改变行列式的值，由引理即得结论. 另一结论类似可证. $\square$

\noindent\textbf{注}\quad 当 $A$ 和 $D$ 都是可逆阵时，我们得到等式
\[
|D||A-BD^{-1}C|=|A||D-CA^{-1}B|.
\]
这个等式称为行列式的降阶公式. 因为当 $D$ 和 $A$ 的阶不等时，可以利用它把高阶行列式的计算化为低阶行列式的计算.

\noindent\textbf{例 2.6.4}\quad 计算下列矩阵的行列式的值:
\[
M=\left(\begin{array}{cccc}
a_1^2 & a_1a_2+1 & \cdots & a_1a_n+1\\
a_2a_1+1 & a_2^2 & \cdots & a_2a_n+1\\
\vdots & \vdots & & \vdots\\
a_na_1+1 & a_na_2+1 & \cdots & a_n^2
\end{array}\right).
\]

\noindent\textbf{解}\quad 将 $M$ 化为
\[
M=-I_n+
\left(\begin{array}{cc}
a_1 & 1\\
a_2 & 1\\
\vdots & \vdots\\
a_n & 1
\end{array}\right)
I_2^{-1}
\left(\begin{array}{cccc}
a_1 & a_2 & \cdots & a_n\\
1 & 1 & \cdots & 1
\end{array}\right).
\]
由降阶公式得到
\[
\begin{aligned}
|M|
&=|I_2|^{-1}|-I_n|\left|I_2+
\left(\begin{array}{cccc}
a_1 & a_2 & \cdots & a_n\\
1 & 1 & \cdots & 1
\end{array}\right)(-I_n)^{-1}
\left(\begin{array}{cc}
a_1 & 1\\
a_2 & 1\\
\vdots & \vdots\\
a_n & 1
\end{array}\right)\right|\\
&=(-1)^n\left|I_2-
\left(\begin{array}{cc}
\sum\limits_{i=1}^n a_i^2 & \sum\limits_{i=1}^n a_i\\
\sum\limits_{i=1}^n a_i & n
\end{array}\right)\right|\\
&=(-1)^n\left((1-n)\left(1-\sum\limits_{i=1}^n a_i^2\right)-\left(\sum\limits_{i=1}^n a_i\right)^2\right).
\end{aligned}
\]

\noindent\textbf{例 2.6.5}\quad 已知 $A$ 和 $D$ 是可逆阵，求下列分块矩阵的逆阵:
\[
\left(\begin{array}{cc}
A & B\\
O & D
\end{array}\right).
\]

\noindent\textbf{解}\quad 设 $A,D$ 分别是 $m,n$ 阶矩阵，对下列分块矩阵进行初等变换，先将第二分块行左乘以 $-BD^{-1}$ 加到第一分块行上去:
\[
\left(\begin{array}{cc|cc}
A & B & I_m & O\\
O & D & O & I_n
\end{array}\right)\longrightarrow
\left(\begin{array}{cc|cc}
A & O & I_m & -BD^{-1}\\
O & D & O & I_n
\end{array}\right),
\]
再用 $A^{-1}$ 和 $D^{-1}$ 分别左乘以第一分块行及第二分块行得到
\[
\left(\begin{array}{cc|cc}
I_m & O & A^{-1} & -A^{-1}BD^{-1}\\
O & I_n & O & D^{-1}
\end{array}\right).
\]
因此原矩阵的逆阵为
\[
\left(\begin{array}{cc}
A^{-1} & -A^{-1}BD^{-1}\\
O & D^{-1}
\end{array}\right).
\]

\subsection*{习题 2.6}

1. 计算下列分块矩阵的乘积:
\[
\begin{array}{ll}
(1)\quad
\left(\begin{array}{cc|ccc}
1 & 0 & 1 & 2 & -1\\
0 & 1 & 3 & 2 & -2\\ \hline
-1 & 4 & 0 & 0 & 0\\
0 & 2 & 0 & 0 & 0
\end{array}\right)
\left(\begin{array}{cc|cc}
2 & -3 & 0 & 0\\
0 & -2 & 0 & 0\\ \hline
1 & 0 & 5 & -1\\
1 & 1 & 0 & 2\\
0 & 0 & 3 & 0
\end{array}\right);\\[3em]
(2)\quad
\left(\begin{array}{cc|cc}
1 & 0 & 0 & 1\\
0 & 1 & 1 & 0\\ \hline
0 & -1 & 1 & 0\\
-1 & 0 & 0 & 1
\end{array}\right)
\left(\begin{array}{cc|cc}
0 & 1 & 0 & -1\\
0 & 0 & -1 & 0\\ \hline
0 & 0 & 0 & 1\\
1 & 0 & 0 & 1
\end{array}\right).
\end{array}
\]

2. 计算下列分块矩阵的乘积，假设其中相乘的分块矩阵都符合相乘的条件:
\[
(1)\quad
\left(\begin{array}{cc}
A_1 & O\\
O & A_2
\end{array}\right)
\left(\begin{array}{cc}
B_{11} & B_{12}\\
B_{21} & B_{22}
\end{array}\right);
\quad
(2)\quad
\left(\begin{array}{ccc}
A_{11} & A_{12} & A_{13}\\
O & A_{22} & A_{23}\\
O & O & A_{33}
\end{array}\right)
\left(\begin{array}{ccc}
B_{11} & B_{12} & B_{13}\\
O & B_{22} & B_{23}\\
O & O & B_{33}
\end{array}\right).
\]

3. 设
\[
A=\left(\begin{array}{cccccc}
0 & 1 & 0 & \cdots & 0 & 0\\
0 & 0 & 1 & \cdots & 0 & 0\\
\vdots & \vdots & \vdots & & \vdots & \vdots\\
0 & 0 & 0 & \cdots & 0 & 1\\
1 & 0 & 0 & \cdots & 0 & 0
\end{array}\right),
\]
求证:
\[
A^k=\left(\begin{array}{cc}
O & I_{n-k}\\
I_k & O
\end{array}\right)\quad (k=1,2,\cdots,n).
\]

4. 设
\[
A=\left(\begin{array}{cccccc}
0 & 1 & 0 & \cdots & 0 & 0\\
0 & 0 & 1 & \cdots & 0 & 0\\
\vdots & \vdots & \vdots & & \vdots & \vdots\\
0 & 0 & 0 & \cdots & 0 & 1\\
0 & 0 & 0 & \cdots & 0 & 0
\end{array}\right),
\]
求证:
\[
A^k=\left(\begin{array}{cc}
O & I_{n-k}\\
O & O
\end{array}\right)\quad (k=1,2,\cdots,n).
\]

5. 求证：分块矩阵的乘法所得的结果与作为普通矩阵的乘法所得的结果是一致的.

6. 求证：分块矩阵的第三类初等变换是普通矩阵的若干次第三类初等变换的复合.

7. 设有分块矩阵 $C=\left(\begin{array}{cc}O&A\\ B&O\end{array}\right)$，其中 $A,B$ 为可逆阵，求 $C$ 的逆阵.

8. 设 $A,B$ 为 $n$ 阶方阵，求证:
\[
\left|\begin{array}{cc}
A & B\\
B & A
\end{array}\right|=|A+B||A-B|.
\]

9. 设 $A,B$ 为 $n$ 阶复方阵，求证:
\[
\left|\begin{array}{cc}
A & -B\\
B & A
\end{array}\right|=|A+\mathrm{i}B||A-\mathrm{i}B|.
\]

10. 设 $A,B$ 为 $n$ 阶方阵且 $AB=BA$，求证:
\[
\left|\begin{array}{cc}
A & -B\\
B & A
\end{array}\right|=|A^2+B^2|.
\]

11. 试用分块矩阵的方法求例 1.5.10 和例 2.5.2 中行列式的值.

12. 求下列矩阵 $A$ 的行列式的值:
\[
(1)\quad A=\left(\begin{array}{ccccc}
0 & 2 & 3 & \cdots & n\\
1 & 0 & 3 & \cdots & n\\
1 & 2 & 0 & \cdots & n\\
\vdots & \vdots & \vdots & & \vdots\\
1 & 2 & 3 & \cdots & 0
\end{array}\right);
\quad
(2)\quad A=\left(\begin{array}{cccc}
1+a_1^2 & a_1a_2 & \cdots & a_1a_n\\
a_2a_1 & 1+a_2^2 & \cdots & a_2a_n\\
\vdots & \vdots & & \vdots\\
a_na_1 & a_na_2 & \cdots & 1+a_n^2
\end{array}\right).
\]

\section*{* §2.7\quad Cauchy-Binet 公式}

我们在本章第五节中，证明了方阵乘积的行列式等于各方阵行列式之积. 现在的问题是：如果 $A$ 是 $m\times n$ 矩阵，$B$ 是 $n\times m$ 矩阵，$AB$ 是 $m$ 阶方阵，则行列式 $|AB|$ 应该等于什么？Cauchy-Binet (柯西--毕内) 公式回答了这个问题. 它可以看成是矩阵乘法的行列式定理的推广.

\noindent\textbf{定理 2.7.1 (Cauchy-Binet 公式)}\quad 设 $A=(a_{ij})$ 是 $m\times n$ 矩阵，$B=(b_{ij})$ 是 $n\times m$ 矩阵. $A\left(\begin{array}{ccc}i_1&\cdots&i_s\\ j_1&\cdots&j_s\end{array}\right)$ 表示 $A$ 的一个 $s$ 阶子式，它由 $A$ 的第 $i_1,\cdots,i_s$ 行与第 $j_1,\cdots,j_s$ 列交点上的元素按原次序排列组成的行列式. 同理可定义 $B$ 的 $s$ 阶子式.

(1) 若 $m>n$，则必有 $|AB|=0$;

(2) 若 $m\leq n$，则必有
\[
|AB|=\sum_{1\leq j_1<j_2<\cdots<j_m\leq n}
A\left(\begin{array}{cccc}
1 & 2 & \cdots & m\\
j_1 & j_2 & \cdots & j_m
\end{array}\right)
B\left(\begin{array}{cccc}
j_1 & j_2 & \cdots & j_m\\
1 & 2 & \cdots & m
\end{array}\right).
\]

\noindent\textbf{证明}\quad 令
\[
C=\left(\begin{array}{cc}
A & O\\
-I_n & B
\end{array}\right).
\]
我们将用不同的方法来计算行列式 $|C|$.

首先，对 $C$ 进行第三类分块初等变换得到矩阵 $M=\left(\begin{array}{cc}O&AB\\ -I_n&B\end{array}\right)$. 事实上，
\[
M=\left(\begin{array}{cc}
I_m & A\\
O & I_n
\end{array}\right)C,
\]
因此 $|M|=|C|$. 用 Laplace 定理来计算 $|M|$，按前 $m$ 行展开得
\[
|M|=(-1)^{(n+1+n+2+\cdots+n+m)+(1+2+\cdots+m)}|-I_n||AB|=(-1)^{n(m+1)}|AB|.
\]
再来计算 $|C|$，用 Laplace 定理按前 $m$ 行展开. 这时若 $m>n$，则前 $m$ 行中任意一个 $m$ 阶子式都至少有一列全为零，因此行列式值等于零，即 $|AB|=0$. 若 $m\leq n$，则由 Laplace 定理得
\[
|C|=\sum_{1\leq j_1<j_2<\cdots<j_m\leq n}
A\left(\begin{array}{cccc}
1 & 2 & \cdots & m\\
j_1 & j_2 & \cdots & j_m
\end{array}\right)
\widehat{C}\left(\begin{array}{cccc}
1 & 2 & \cdots & m\\
j_1 & j_2 & \cdots & j_m
\end{array}\right),
\]
其中 $\widehat{C}\left(\begin{array}{cccc}1&2&\cdots&m\\ j_1&j_2&\cdots&j_m\end{array}\right)$ 是 $A\left(\begin{array}{cccc}1&2&\cdots&m\\ j_1&j_2&\cdots&j_m\end{array}\right)$ 在矩阵 $C$ 中的代数余子式. 显然
\[
\widehat{C}\left(\begin{array}{cccc}
1 & 2 & \cdots & m\\
j_1 & j_2 & \cdots & j_m
\end{array}\right)
=(-1)^{\frac12 m(m+1)+(j_1+j_2+\cdots+j_m)}
|-e_{i_1},-e_{i_2},\cdots,-e_{i_{n-m}},B|,
\]
其中 $i_1,i_2,\cdots,i_{n-m}$ 是 $C$ 中前 $n$ 列去掉 $j_1,j_2,\cdots,j_m$ 列后余下的列序数. $e_{i_1},e_{i_2},\cdots,e_{i_{n-m}}$ 是相应的 $n$ 维标准单位列向量 (标准单位向量定义参见本章复习题 1). 记
\[
|N|=|-e_{i_1},-e_{i_2},\cdots,-e_{i_{n-m}},B|,
\]
现来计算 $|N|$. $|N|$ 用 Laplace 定理按前 $n-m$ 列展开. 注意只有一个子式非零，其值等于 $|-I_{n-m}|=(-1)^{n-m}$. 而这个子式的余子式为
\[
B\left(\begin{array}{cccc}
j_1 & j_2 & \cdots & j_m\\
1 & 2 & \cdots & m
\end{array}\right).
\]
因此
\[
|N|=(-1)^{(n-m)+(i_1+i_2+\cdots+i_{n-m})+(1+2+\cdots+n-m)}
B\left(\begin{array}{cccc}
j_1 & j_2 & \cdots & j_m\\
1 & 2 & \cdots & m
\end{array}\right).
\]
注意到
\[
(i_1+i_2+\cdots+i_{n-m})+(j_1+j_2+\cdots+j_m)=1+2+\cdots+n.
\]
综合上面的结论，通过简单计算不难得到
\[
|AB|=\sum_{1\leq j_1<j_2<\cdots<j_m\leq n}
A\left(\begin{array}{cccc}
1 & 2 & \cdots & m\\
j_1 & j_2 & \cdots & j_m
\end{array}\right)
B\left(\begin{array}{cccc}
j_1 & j_2 & \cdots & j_m\\
1 & 2 & \cdots & m
\end{array}\right). \quad \square
\]

下面的定理是 Cauchy-Binet 公式的进一步推广，它告诉我们如何求矩阵乘积的 $r$ 阶子式.

\noindent\textbf{定理 2.7.2}\quad 设 $A=(a_{ij})$ 是 $m\times n$ 矩阵，$B=(b_{ij})$ 是 $n\times m$ 矩阵，$r$ 是一个正整数且 $r\leq m$.

(1) 若 $r>n$，则 $AB$ 的任意一个 $r$ 阶子式等于零;

(2) 若 $r\leq n$，则 $AB$ 的 $r$ 阶子式
\[
AB\left(\begin{array}{cccc}
i_1 & i_2 & \cdots & i_r\\
j_1 & j_2 & \cdots & j_r
\end{array}\right)
\]
\[
=\sum_{1\leq k_1<k_2<\cdots<k_r\leq n}
A\left(\begin{array}{cccc}
i_1 & i_2 & \cdots & i_r\\
k_1 & k_2 & \cdots & k_r
\end{array}\right)
B\left(\begin{array}{cccc}
k_1 & k_2 & \cdots & k_r\\
j_1 & j_2 & \cdots & j_r
\end{array}\right).
\]

\noindent\textbf{证明}\quad 设 $C=AB$，则 $C=(c_{ij})$ 是 $m$ 阶矩阵且
\[
c_{ij}=a_{i1}b_{1j}+a_{i2}b_{2j}+\cdots+a_{in}b_{nj}.
\]
因此
\[
C\left(\begin{array}{cccc}
i_1 & i_2 & \cdots & i_r\\
j_1 & j_2 & \cdots & j_r
\end{array}\right)=
\left|\begin{array}{cccc}
a_{i_11} & a_{i_12} & \cdots & a_{i_1n}\\
a_{i_21} & a_{i_22} & \cdots & a_{i_2n}\\
\vdots & \vdots & & \vdots\\
a_{i_r1} & a_{i_r2} & \cdots & a_{i_rn}
\end{array}\right|
\left|\begin{array}{cccc}
b_{1j_1} & b_{1j_2} & \cdots & b_{1j_r}\\
b_{2j_1} & b_{2j_2} & \cdots & b_{2j_r}\\
\vdots & \vdots & & \vdots\\
b_{nj_1} & b_{nj_2} & \cdots & b_{nj_r}
\end{array}\right|.
\]
由定理 2.7.1 可知：当 $r>n$ 时，
\[
C\left(\begin{array}{cccc}
i_1 & i_2 & \cdots & i_r\\
j_1 & j_2 & \cdots & j_r
\end{array}\right)=0;
\]
当 $r\leq n$ 时，
\[
C\left(\begin{array}{cccc}
i_1 & i_2 & \cdots & i_r\\
j_1 & j_2 & \cdots & j_r
\end{array}\right)
=\sum_{1\leq k_1<k_2<\cdots<k_r\leq n}
A\left(\begin{array}{cccc}
i_1 & i_2 & \cdots & i_r\\
k_1 & k_2 & \cdots & k_r
\end{array}\right)
B\left(\begin{array}{cccc}
k_1 & k_2 & \cdots & k_r\\
j_1 & j_2 & \cdots & j_r
\end{array}\right). \quad \square
\]

矩阵 $A$ 的子式
\[
A\left(\begin{array}{cccc}
i_1 & i_2 & \cdots & i_r\\
j_1 & j_2 & \cdots & j_r
\end{array}\right)
\]
如果满足条件 $i_1=j_1,i_2=j_2,\cdots,i_r=j_r$，则称为主子式.

\noindent\textbf{推论 2.7.1}\quad 设 $A$ 是 $m\times n$ 实矩阵，则矩阵 $AA'$ 的任一主子式都非负.

\noindent\textbf{证明}\quad 若 $r\leq n$，则由定理 2.7.2 得到
\[
AA'\left(\begin{array}{cccc}
i_1 & i_2 & \cdots & i_r\\
i_1 & i_2 & \cdots & i_r
\end{array}\right)=
\sum_{1\leq k_1<k_2<\cdots<k_r\leq n}
A\left(\begin{array}{cccc}
i_1 & i_2 & \cdots & i_r\\
k_1 & k_2 & \cdots & k_r
\end{array}\right)^2\geq 0;
\]
若 $r>n$，则 $AA'$ 的任一 $r$ 阶主子式都等于零，结论也成立. $\square$

下面介绍 Cauchy-Binet 公式的两个重要应用. 它们分别是著名的 Lagrange (拉格朗日) 恒等式和 Cauchy-Schwarz (柯西--许瓦兹) 不等式. 这两个结论也可以用其他方法证明，但用矩阵方法显得非常简洁.

\noindent\textbf{例 2.7.1}\quad 证明 Lagrange 恒等式 $(n\geq 2)$:
\[
\left(\sum_{i=1}^n a_i^2\right)\left(\sum_{i=1}^n b_i^2\right)-\left(\sum_{i=1}^n a_ib_i\right)^2
=\sum_{1\leq i<j\leq n}(a_ib_j-a_jb_i)^2.
\]

\noindent\textbf{证明}\quad 左边的式子等于
\[
\left|\begin{array}{cc}
\sum\limits_{i=1}^n a_i^2 & \sum\limits_{i=1}^n a_ib_i\\
\sum\limits_{i=1}^n a_ib_i & \sum\limits_{i=1}^n b_i^2
\end{array}\right|,
\]
这个行列式对应的矩阵可化为
\[
\left(\begin{array}{cccc}
a_1 & a_2 & \cdots & a_n\\
b_1 & b_2 & \cdots & b_n
\end{array}\right)
\left(\begin{array}{cc}
a_1 & b_1\\
a_2 & b_2\\
\vdots & \vdots\\
a_n & b_n
\end{array}\right).
\]
用 Cauchy-Binet 公式得
\[
\left|\begin{array}{cc}
\sum\limits_{i=1}^n a_i^2 & \sum\limits_{i=1}^n a_ib_i\\
\sum\limits_{i=1}^n a_ib_i & \sum\limits_{i=1}^n b_i^2
\end{array}\right|
=\sum_{1\leq i<j\leq n}
\left|\begin{array}{cc}
a_i & a_j\\
b_i & b_j
\end{array}\right|
\left|\begin{array}{cc}
a_i & b_i\\
a_j & b_j
\end{array}\right|
=\sum_{1\leq i<j\leq n}(a_ib_j-a_jb_i)^2. \quad \square
\]

\noindent\textbf{例 2.7.2}\quad 设 $a_i,b_i$ 都是实数，证明 Cauchy-Schwarz 不等式:
\[
\left(\sum_{i=1}^n a_i^2\right)\left(\sum_{i=1}^n b_i^2\right)\geq
\left(\sum_{i=1}^n a_ib_i\right)^2.
\]

\noindent\textbf{证明}\quad 由例 2.7.1，恒等式右边总非负，即得结论. $\square$

\subsection*{习题 2.7}

1. 设 $n\geq 3$，求证下列行列式的值为零:
\[
|A|=\left|\begin{array}{cccc}
1+x_1y_1 & 1+x_1y_2 & \cdots & 1+x_1y_n\\
1+x_2y_1 & 1+x_2y_2 & \cdots & 1+x_2y_n\\
\vdots & \vdots & & \vdots\\
1+x_ny_1 & 1+x_ny_2 & \cdots & 1+x_ny_n
\end{array}\right|.
\]

2. 设 $A$ 是 $n$ 阶实方阵且 $AA'=I_n$. 求证：若 $1\leq i_1<i_2<\cdots<i_k\leq n$，则
\[
\sum_{1\leq j_1<j_2<\cdots<j_k\leq n}
A\left(\begin{array}{cccc}
i_1 & i_2 & \cdots & i_k\\
j_1 & j_2 & \cdots & j_k
\end{array}\right)^2=1.
\]

3. 设 $A,B$ 是 $n$ 阶方阵，$r$ 是小于等于 $n$ 的正整数，求证：$AB$ 和 $BA$ 的所有 $r$ 阶主子式之和相等.

4. 设 $A,B$ 都是 $m\times n$ 实矩阵，求证:
\[
|AA'||BB'|\geq |AB'|^2.
\]

5. 设 $A$ 是 $m\times n$ 复矩阵，求证：矩阵 $A\overline{A}'$ 的任一主子式都非负.

6. 证明复数形式的 Cauchy-Schwarz 不等式:
\[
\left(\sum_{i=1}^n |a_i|^2\right)\left(\sum_{i=1}^n |b_i|^2\right)\geq
\left|\sum_{i=1}^n a_i\overline{b_i}\right|^2,
\]
其中 $a_i,b_i$ 都是复数，$n$ 是大于等于 2 的正整数.

7. 将习题 4 推广到复数形式并证明之.

\section*{历史与展望}

矩阵是一个重要的数学研究对象，它伴随着线性方程组、线性变换、二次型和双线性型而自然出现，在数论、几何和分析等数学分支，以及物理和工程等学科中均有着广泛的应用. 简单来说，矩阵就是一个长方形数组或二维数组. 矩阵的雏形最早出现在公元前 200 年的著作《九章算术》中，根据该著作所述可知，那时的中国人就利用矩阵和消元法来求解三元线性方程组.

与行列式一样，矩阵也是解决线性方程组求解问题的产物. 由第一章的介绍可知，利用行列式求解线性方程组的研究起源于 Leibniz. 到了 18 世纪，线性方程组的研究通常被归类在行列式之下，因此并没有研究方程个数与未定元个数不相等的线性方程组.

1800 年左右 Gauss 给出了利用矩阵求解线性方程组的系统方法，现称之为 Gauss 消元法. 这一方法处理了方程个数与未定元个数不相等的线性方程组，尽管当时 Gauss 并没有使用矩阵的符号.

Gauss 在著作《算术研究》中研究了二元二次型的算术理论. 虽然 Gauss 并没有明确地给出矩阵的术语，但他将变量之间的线性变换含蓄地表示为矩阵，将线性变换的复合含蓄地定义为矩阵的乘积. 当然，他仅仅是考虑 $2\times 2$ 矩阵和 $3\times 3$ 矩阵.

1844 年 Eisenstein (爱森斯坦因) 给出了矩阵乘法的概念，用于化简线性方程组求解过程中的变量代换. 1850 年 Sylvester 创造了矩阵的术语.

Cayley 在 1850 年和 1858 年的两篇论文中，正式地给出了 $m\times n$ 矩阵的概念；他指出“矩阵把自己组合成单个实体”，并认识到矩阵在简化线性方程组的求解和线性变换的复合上的用处；他给出了矩阵加法、乘法和数乘的定义，给出了单位阵和方阵的逆阵的概念；说明了在某种条件下，如何用逆阵来求解 $n$ 个未定元 $n$ 个方程的线性方程组.

在 1858 年的论文《矩阵理论的研究报告》中，Cayley 证明了 Cayley-Hamilton (凯莱--哈密顿) 定理，即方阵适合它自己的特征多项式. 这个证明由他对 $2\times 2$ 矩阵和 $3\times 3$ 矩阵的计算组成. 他指出这个结果可推广到更高阶，但他补充道：“我认为没有必要花这个力气去正式证明这个定理对任意阶矩阵都成立.” Hamilton 在他的四元数研究中独立地证明了这个定理 ($n=4$ 的情形，但没有使用矩阵的符号). Cayley 在另一篇文章中使用矩阵解决了一个非常有意义的问题，即 Cayley-Hermite 问题，它要求确定使得 $n$ 元二次型保持不变的所有的线性变换.

作为矩阵理论的奠基者，Cayley 极大地增强了将矩阵视为一种符号代数的重要思想. 特别地，他使用了单个字母来表示矩阵，这在矩阵代数的发展史上是具有重大意义的一步. 但是他写于 1850 年代的论文直到 1880 年代为止，在英国以外并未被重视.

与行列式理论的研究在 20 世纪初就偃旗息鼓不同，矩阵理论的研究直到今天仍然方兴未艾. 随着大数据科学、人工智能和 5G 技术等高新科技的异军突起，矩阵理论的研究将越来越重要，对现代科技的影响也越发深远. 矩阵理论的进一步发展及其应用将会在第六章、第七章、第八章和第九章继续介绍.

\section*{复习题二}

1. $n$ 维标准单位列向量是指下面 $n$ 个 $n$ 维列向量:
\[
e_1=\left(\begin{array}{c}1\\0\\ \vdots\\0\end{array}\right),\quad
e_2=\left(\begin{array}{c}0\\1\\ \vdots\\0\end{array}\right),\quad \cdots,\quad
e_n=\left(\begin{array}{c}0\\0\\ \vdots\\1\end{array}\right).
\]
向量组 $e_1',e_2',\cdots,e_n'$ 则被称为 $n$ 维标准单位行向量. 设 $f_1,f_2,\cdots,f_m$ 是 $m$ 维标准单位列向量. 证明标准单位向量有下列基本性质:

(1) 若 $i\neq j$，则 $e_i'e_j=0$，而 $e_i'e_i=1$;

(2) 若 $A=(a_{ij})$ 是 $m\times n$ 矩阵，则 $Ae_i$ 是 $A$ 的第 $i$ 个列向量，$f_i'A$ 是 $A$ 的第 $i$ 个行向量;

(3) 若 $A=(a_{ij})$ 是 $m\times n$ 矩阵，则 $f_i'Ae_j=a_{ij}$;

(4) 设 $A,B$ 都是 $m\times n$ 矩阵，则 $A=B$ 当且仅当 $Ae_i=Be_i(1\leq i\leq n)$，也当且仅当 $f_i'A=f_i'B(1\leq i\leq m)$.

2. $n$ 阶基础矩阵 (又称初级矩阵) 是指 $n^2$ 个 $n$ 阶矩阵 $\{E_{ij},1\leq i,j\leq n\}$. 这里 $E_{ij}$ 是一个 $n$ 阶矩阵，它的第 $(i,j)$ 元素等于 1，其他元素全为 0. 基础矩阵也可以看成是标准单位向量的积: $E_{ij}=e_ie_j'$. 证明基础矩阵的下列性质:

(1) 若 $j\neq k$，则 $E_{ij}E_{kl}=O$;

(2) 若 $j=k$，则 $E_{ij}E_{kl}=E_{il}$;

(3) 若 $A$ 是 $n$ 阶矩阵且 $A=(a_{ij})$，则 $A=\sum\limits_{i,j=1}^n a_{ij}E_{ij}$;

(4) 若 $A$ 是 $n$ 阶矩阵且 $A=(a_{ij})$，则 $E_{ij}A$ 的第 $i$ 行是 $A$ 的第 $j$ 行，$E_{ij}A$ 的其他行全为零;

(5) 若 $A$ 是 $n$ 阶矩阵且 $A=(a_{ij})$，则 $AE_{ij}$ 的第 $j$ 列是 $A$ 的第 $i$ 列，$AE_{ij}$ 的其他列全为零;

(6) 若 $A$ 是 $n$ 阶矩阵且 $A=(a_{ij})$，则 $E_{ij}AE_{kl}=a_{jk}E_{il}$.

3. 设 $A$ 为 $n$ 阶方阵，求证：$A$ 是反对称阵的充分必要条件是对任意的 $n$ 维列向量 $\alpha$，有
\[
\alpha'A\alpha=0.
\]

4. 设 $A$ 是 $n$ 阶上三角阵且主对角线上元素全为零，求证：$A^n=O$.

5. 若 $A,B$ 都是由非负实数组成的矩阵且 $AB$ 有一行全为零，求证：或者 $A$ 有一行全为零，或者 $B$ 有一列全为零.

6. 设 $A$ 是二阶矩阵，若存在 $n>2$，使得 $A^n=O$，求证：$A^2=O$.

7. 下列形状的矩阵称为循环矩阵:
\[
\left(\begin{array}{ccccc}
a_1 & a_2 & a_3 & \cdots & a_n\\
a_n & a_1 & a_2 & \cdots & a_{n-1}\\
a_{n-1} & a_n & a_1 & \cdots & a_{n-2}\\
\vdots & \vdots & \vdots & & \vdots\\
a_2 & a_3 & a_4 & \cdots & a_1
\end{array}\right).
\]
求证：同阶循环矩阵之积仍是循环矩阵.

8. 设 $n$ 阶方阵 $A$ 的每一行元素之和等于常数 $c$，求证:

(1) 对任意的正整数 $k$，$A^k$ 的每一行元素之和等于常数 $c^k$;

(2) 若 $A$ 为可逆阵，则 $c\neq 0$ 并且 $A^{-1}$ 的每一行元素之和等于 $c^{-1}$.

9. 设 $A$ 是奇数阶矩阵，满足 $AA'=I_n$ 且 $|A|>0$，求证：$I_n-A$ 是奇异阵.

10. 设 $A,B$ 为 $n$ 阶可逆阵，满足 $A^2=B^2$ 且 $|A|+|B|=0$，求证：$A+B$ 是奇异阵.

11. 设 $A,B,AB-I_n$ 都是 $n$ 阶可逆阵，证明：$A-B^{-1}$ 与 $(A-B^{-1})^{-1}-A^{-1}$ 均可逆，并求它们的逆阵.

12. 设 $A$ 为 $m\times n$ 矩阵，$B$ 为 $n\times m$ 矩阵，使得 $I_m+AB$ 可逆，求证：$I_n+BA$ 也可逆.

13. 设 $A,B,A-B$ 都是 $n$ 阶可逆阵，证明:
\[
B^{-1}-A^{-1}=\left(B+B(A-B)^{-1}B\right)^{-1}.
\]

14. 设 $A$ 是 $n$ 阶可逆阵，$\alpha,\beta$ 是 $n$ 维列向量，且 $1+\beta'A^{-1}\alpha\neq 0$. 求证:
\[
(A+\alpha\beta')^{-1}=A^{-1}-\frac{1}{1+\beta'A^{-1}\alpha}A^{-1}\alpha\beta'A^{-1}.
\]
注：上述公式称为 Sherman-Morrison (谢尔曼--莫里森) 公式.

15. 求证：$n$ 阶方阵 $A$ 是奇异阵的充分必要条件是存在非零的 $n$ 维列向量 $x$，使得 $Ax=0$.

16. 设 $A$ 为 $n$ 阶实反对称阵，证明：$I_n-A$ 是非异阵.

17. 设 $A$ 为 $n$ 阶可逆阵，求证：只用第三类初等变换就可以将 $A$ 化为如下形状:
\[
\operatorname{diag}\{1,1,\cdots,|A|\}.
\]

18. 求证：任一 $n$ 阶矩阵均可表示为形如 $I_n+a_{ij}E_{ij}$ 这样的矩阵之积，其中 $E_{ij}$ 是 $n$ 阶基础矩阵.

19. 求下列 $n$ 阶矩阵的逆阵，其中 $a_i\neq 0(1\leq i\leq n)$:
\[
A=\left(\begin{array}{ccccc}
1+a_1 & 1 & 1 & \cdots & 1\\
1 & 1+a_2 & 1 & \cdots & 1\\
1 & 1 & 1+a_3 & \cdots & 1\\
\vdots & \vdots & \vdots & & \vdots\\
1 & 1 & 1 & \cdots & 1+a_n
\end{array}\right).
\]

20. 设 $A,B$ 为 $n$ 阶矩阵，求证：$(AB)^*=B^*A^*$.

21. 设 $A$ 为 $n$ 阶矩阵，求证：$|A^*|=|A|^{n-1}$.

22. 设 $A$ 为 $n(n>2)$ 阶矩阵，求证：$(A^*)^*=|A|^{n-2}A$.

23. 设 $A$ 为 $m$ 阶矩阵，$B$ 为 $n$ 阶矩阵，求分块对角阵 $C$ 的伴随矩阵:
\[
C=\left(\begin{array}{cc}
A & O\\
O & B
\end{array}\right).
\]

24. 已知
\[
A^*=\left(\begin{array}{ccc}
1 & -2 & 1\\
0 & 2 & -2\\
-1 & 2 & 1
\end{array}\right),
\]
试求 $A$.

25. 设 $n$ 阶矩阵
\[
A=\left(\begin{array}{ccccc}
2 & 2 & 2 & \cdots & 2\\
0 & 1 & 1 & \cdots & 1\\
0 & 0 & 1 & \cdots & 1\\
\vdots & \vdots & \vdots & & \vdots\\
0 & 0 & 0 & \cdots & 1
\end{array}\right),
\]
试求 $\sum\limits_{i,j=1}^n A_{ij}$，即 $|A|$ 的所有代数余子式之和.

26. 设 $A=(a_{ij})$ 是 $n$ 阶矩阵，则 $A$ 主对角线上元素之和
\[
a_{11}+a_{22}+\cdots+a_{nn}
\]
称为矩阵 $A$ 的迹，记为 $\operatorname{tr}(A)$. 设 $A,B$ 是 $n$ 阶矩阵，$k$ 是常数，求证:
\[
(1)\ \operatorname{tr}(A+B)=\operatorname{tr}(A)+\operatorname{tr}(B);\qquad
(2)\ \operatorname{tr}(kA)=k\operatorname{tr}(A);
\]
\[
(3)\ \operatorname{tr}(A')=\operatorname{tr}(A);\qquad
(4)\ \operatorname{tr}(AB)=\operatorname{tr}(BA).
\]

27. 求证：不存在 $n$ 阶矩阵 $A,B$，使得 $AB-BA=kI_n(k\neq 0)$.

28. 证明下列结论:

(1) 若 $A$ 是 $m\times n$ 实矩阵，则 $\operatorname{tr}(AA')\geq 0$，等号成立的充分必要条件是 $A=O$;

(2) 若 $A$ 是 $m\times n$ 复矩阵，则 $\operatorname{tr}(A\overline{A}')\geq 0$，等号成立的充分必要条件是 $A=O$.

29. 设 $A_i(i=1,2,\cdots,k)$ 是实对称阵且 $A_1^2+A_2^2+\cdots+A_k^2=O$，证明：每个 $A_i=O$.

30. 证明下列结论:

(1) 设 $n$ 阶实矩阵 $A$ 适合 $A'=-A$，如果存在同阶实矩阵 $B$，使得 $AB=B$，则 $B=O$;

(2) 设 $n$ 阶复矩阵 $A$ 适合 $\overline{A}'=-A$，如果存在同阶复矩阵 $B$，使得 $AB=B$，则 $B=O$.

31. 设 $A$ 为 $n$ 阶实矩阵，求证：$\operatorname{tr}(A^2)\leq \operatorname{tr}(AA')$，等号成立当且仅当 $A$ 是对称阵.

32. 设 $A,B$ 是 $n$ 阶矩阵，使得 $\operatorname{tr}(ABC)=\operatorname{tr}(CBA)$ 对任意 $n$ 阶矩阵 $C$ 成立，求证：$AB=BA$.

33. 设 $f$ 是数域 $\mathbb{F}$ 上 $n$ 阶矩阵集合到 $\mathbb{F}$ 的一个映射，它满足下列条件:

(1) 对任意的 $n$ 阶矩阵 $A,B$，$f(A+B)=f(A)+f(B)$;

(2) 对任意的 $n$ 阶矩阵 $A$ 和 $\mathbb{F}$ 中数 $k$，$f(kA)=kf(A)$;

(3) 对任意的 $n$ 阶矩阵 $A,B$，$f(AB)=f(BA)$;

(4) $f(I_n)=n$.

求证：$f$ 就是迹，即 $f(A)=\operatorname{tr}(A)$ 对一切 $\mathbb{F}$ 上 $n$ 阶矩阵 $A$ 成立.

34. 计算矩阵 $A$ 的行列式的值:
\[
A=\left(\begin{array}{ccccc}
\cos\theta & \cos2\theta & \cos3\theta & \cdots & \cos n\theta\\
\cos n\theta & \cos\theta & \cos2\theta & \cdots & \cos(n-1)\theta\\
\cos(n-1)\theta & \cos n\theta & \cos\theta & \cdots & \cos(n-2)\theta\\
\vdots & \vdots & \vdots & & \vdots\\
\cos2\theta & \cos3\theta & \cos4\theta & \cdots & \cos\theta
\end{array}\right).
\]

35. 计算矩阵 $A$ 的行列式的值:
\[
A=\left(\begin{array}{ccccc}
a_1 & a_2 & a_3 & \cdots & a_n\\
-a_n & a_1 & a_2 & \cdots & a_{n-1}\\
-a_{n-1} & -a_n & a_1 & \cdots & a_{n-2}\\
\vdots & \vdots & \vdots & & \vdots\\
-a_2 & -a_3 & -a_4 & \cdots & a_1
\end{array}\right).
\]

36. 计算矩阵 $A$ 的行列式的值:
\[
A=\left(\begin{array}{cccc}
a_1-b_1 & a_1-b_2 & \cdots & a_1-b_n\\
a_2-b_1 & a_2-b_2 & \cdots & a_2-b_n\\
\vdots & \vdots & & \vdots\\
a_n-b_1 & a_n-b_2 & \cdots & a_n-b_n
\end{array}\right).
\]

37. 设 $n\geq 3$，证明下列矩阵是奇异阵:
\[
A=\left(\begin{array}{cccc}
\cos(\alpha_1-\beta_1) & \cos(\alpha_1-\beta_2) & \cdots & \cos(\alpha_1-\beta_n)\\
\cos(\alpha_2-\beta_1) & \cos(\alpha_2-\beta_2) & \cdots & \cos(\alpha_2-\beta_n)\\
\vdots & \vdots & & \vdots\\
\cos(\alpha_n-\beta_1) & \cos(\alpha_n-\beta_2) & \cdots & \cos(\alpha_n-\beta_n)
\end{array}\right).
\]

38. 计算矩阵 $A$ 的行列式的值，其中 $a_i\neq 0(1\leq i\leq n)$:
\[
A=\left(\begin{array}{cccc}
0 & a_1+a_2 & \cdots & a_1+a_n\\
a_2+a_1 & 0 & \cdots & a_2+a_n\\
\vdots & \vdots & & \vdots\\
a_n+a_1 & a_n+a_2 & \cdots & 0
\end{array}\right).
\]

39. 设 $A,B,C,D$ 都是 $n$ 阶矩阵，求证:
\[
|M|=\left|\begin{array}{cccc}
A & B & C & D\\
B & A & D & C\\
C & D & A & B\\
D & C & B & A
\end{array}\right|
=|A+B+C+D||A+B-C-D||A-B+C-D||A-B-C+D|.
\]

40. 设 $A,B$ 都是 $n$ 阶矩阵，求证:
\[
|A+B|=|A|+|B|+
\sum_{1\leq k\leq n-1}
\left(
\sum_{\substack{1\leq i_1<i_2<\cdots<i_k\leq n\\
1\leq j_1<j_2<\cdots<j_k\leq n}}
A\left(\begin{array}{cccc}
i_1 & i_2 & \cdots & i_k\\
j_1 & j_2 & \cdots & j_k
\end{array}\right)
\widehat{B}\left(\begin{array}{cccc}
i_1 & i_2 & \cdots & i_k\\
j_1 & j_2 & \cdots & j_k
\end{array}\right)
\right).
\]

41. 设 $f$ 是数域 $\mathbb{F}$ 上 $n$ 阶矩阵集合到 $\mathbb{F}$ 的一个映射，使得对任意的 $n$ 阶方阵 $A$，任意的指标 $1\leq i\leq n$，以及任意的常数 $c\in\mathbb{F}$，满足下列条件:

(1) 设 $A$ 的第 $i$ 列是方阵 $B$ 和 $C$ 的第 $i$ 列之和，且 $A$ 的其余列与 $B$ 和 $C$ 的对应列完全相同，则 $f(A)=f(B)+f(C)$;

(2) 将 $A$ 的第 $i$ 列乘以常数 $c$ 得到方阵 $B$，则 $f(B)=cf(A)$;

(3) 对换 $A$ 的任意两列得到方阵 $B$，则 $f(B)=-f(A)$;

(4) $f(I_n)=1$.

求证：$f(A)=|A|$.

42. 设 $A$ 是一个 $n$ 阶方阵，求证：存在一个正数 $a$，使得对任意的 $0<t<a$，矩阵 $tI_n+A$ 都是非异阵.

43. 设 $A,B,C,D$ 是 $n$ 阶矩阵且 $AC=CA$，求证:
\[
\left|\begin{array}{cc}
A & B\\
C & D
\end{array}\right|=|AD-CB|.
\]

\end{document}
